\documentclass{article}
\usepackage{decomp}
\begin{document}
\decompnote{Collection markers: a cited work that is in the local collection is tagged \localcopy{<key>} after its reference entry; it is decompiled in decompiled/<key>/. A reference to another version of a local work (e.g. its preprint) is tagged \localother{<key>}. A reference without that tag is not in the local collection and needs an external pull. Tags added by \texttt{\detokenize{tools/mark_local.py}}.}

\sourcefile{41592\_2023\_Article\_1866.pdf}


\decompnote{Publisher's version of record (Springer Nature typeset PDF, Nature Methods 20:836--840, 25 pages: 5 pages of article, 6 pages of Methods and back matter, 11 pages of Extended Data Figs.\ 1--7 with captions, and the 3-page Nature Portfolio reporting summary as pp.\,23--25; the reporting summary is identical to the separate file 41592\_2023\_1866\_MOESM2\_ESM.pdf and is decompiled in reporting\_summary.tex, not here). Left out: running headers (``Brief Communication'' + DOI) and footers (journal, volume, pages, page number) on every page, the ``Check for updates'' button, the open-access logo, and the ORCID and e-mail icons after the author names. The journal block is given once below.}

\pdfpage{1}
\noindent nature methods\par
\noindent Brief Communication\quad https://doi.org/10.1038/s41592-023-01866-3\par

\section*{High-contrast en bloc staining of mouse whole-brain and human brain samples for EM-based connectomics}

\noindent Kun Song, Zhihui Feng \& Moritz Helmstaedter

\begin{quote}\small
\noindent Department of Connectomics, Max Planck Institute for Brain Research, Frankfurt, Germany. e-mail: kun.song@brain.mpg.de; mh@brain.mpg.de
\end{quote}
\decompnote{The affiliation and e-mail line is printed at the bottom of the left column of p.\,1; it is placed here after the author list. The e-mail icons mark Kun Song and Moritz Helmstaedter as corresponding authors.}

\noindent Received: 28 February 2022\\
Accepted: 27 March 2023\\
Published online: 8 May 2023

\begin{quote}
Connectomes of human cortical gray matter require high-contrast homogeneously stained samples sized at least 2\,mm on a side, and a mouse whole-brain connectome requires samples sized at least 5–10\,mm on a side. Here we report en bloc staining and embedding protocols for these and other applications, removing a key obstacle for connectomic analyses at the mammalian whole-brain level.
\end{quote}

The dense and homogeneous deposition of heavy metals into brain tissue that leads to high membrane contrast for electron-based imaging is a prerequisite for synaptic-resolution connectomics. Ever since the development of ‘reduced-osmium’ protocols\textsuperscript{1–3}, en bloc staining of tissue samples up to a thickness of about 100–200\,µm was possible. Beyond such sample sizes, however, substantial staining gradients occurred, which limited connectomic analyses to smaller samples\textsuperscript{4,5}. With the development of a modified staining protocol\textsuperscript{6}, samples up to about 1\,mm in size could be homogeneously stained, an important step to allow millimeter-size connectomic data acquisition (Fig. 1a). This protocol has been widely applied since\textsuperscript{7–10}.

However, with the ambition to obtain connectomes from even larger samples, in particular, samples that encompass the gray matter depth of the human cortex (at least 2\,mm in size) and samples corresponding to entire brains of small mammals such as mice\textsuperscript{11} or even humans\textsuperscript{12}, the need for improved protocols became obvious. In spite of the promising initial attempts for mouse whole-brain staining\textsuperscript{13}, there is thus far no reliable protocol for en bloc staining of multi-mm-to-centimeter-scale samples with high staining contrast. This is particularly challenging, as EM imaging and reconstruction have progressed to a stage at which large-sample analyses appear to become possible\textsuperscript{7,8,14}.

Here, we report such protocols for samples 2\,mm in size (Fig. 1a,b, 8\,mm\textsuperscript{3}), mouse hemispheres and whole brains (Fig. 1a,c, 250–500\,mm\textsuperscript{3}; see ref. 15 for an initial version of our protocol) and human cortex samples (Fig. 2c–e, 12\,mm\textsuperscript{3}). Their development required the concomitant solution of the following problems: recurring staining inhomogeneity, sample instability that leads to breakages (especially in hemispheres and whole-brain samples) and inhomogeneous resin infiltration, as described in the following.

For protocol development, we initially used X-ray microtomography (µCT) imaging to assess staining gradients\textsuperscript{13,16}. Additionally, we applied low-vacuum scanning electron microscopy (SEM)\textsuperscript{17,18} to evaluate ultrastructural contrast. This was of particular importance, as samples that had a homogeneous appearance in µCT could reveal insufficient membrane contrast or damaged ultrastructure when analyzed by electron microscopy (EM) (Extended Data Figs. 1 and 2).

To achieve the ultimate goal of a 500-fold enlarged infiltration volume (mouse whole brain versus cubic millimeter, Fig. 1a), we first worked on the staining of a volume about eightfold to tenfold larger (‘2-mm’ samples, Fig. 1a,b). When we applied the available 1-mm protocol\textsuperscript{6} to samples 2\,mm in size (Fig. 1b), we observed strong staining gradients (Fig. 1a and Extended Data Fig. 1a,b) and incomplete resin infiltration (Extended Data Fig. 1j). We then used µCT to investigate at which steps in the protocol the gradients occurred, analyzed their co-dependence and found that, by prolonging relevant incubation steps, targeted separation from washing steps, replacing thiocarbohydrazide (TCH) with pyrogallol (Pg) as proposed in ref. 13 and introducing additional steps to ensure sample stability, we could obtain gradient-free high-contrast samples sized about 2\,mm on a side (Fig. 1b, bottom; see Supplementary Results for a more detailed description of the protocol development and insights from these investigations). For obtaining homogeneous resin infiltration, we considered the fact that the epoxy resin blender would undergo polymerization during infiltration, which would increase the viscosity of the blender, and, once the polymerization process has crossed the gel point, no more diffusion would be possible\textsuperscript{19}. Thus, the practical strategy for improving resin infiltration was to slow down polymerization reactions and maintain low viscosity. To do so, we kept all resin-infiltration steps at 4\,°C\textsuperscript{19,20}, added a step with 95\% resin in 5\% acetone (this small amount of acetone substantially decreases the resin viscosity\textsuperscript{21}) and extended the incubation times (for details, see Supplementary Table 1). By these modifications, we infiltrated both \pdfpage{2}Spurr’s and Epon epoxy resin into the sample center, yielding homogeneously embedded samples 2\,mm in size (Extended Data Fig. 1j,k).

To obtain staining and embedding protocols for mouse hemispheres about 25-fold larger (Extended Data Fig. 3) and adult mouse whole brains about 50-fold larger (Fig. 1c–e), we first extended the

\begin{center}
\includegraphics[width=0.95\linewidth,height=0.55\textheight,keepaspectratio]{figures/fig1.png}
\end{center}
\noindent\textbf{Fig. 1 | High-contrast staining protocol for brain samples 2\,mm to 1 centimeter in size, applied to human cortical tissue and adult mouse whole brains.} \textbf{a}, Overview of en bloc staining protocol development: reduced-osmium-TCH protocols\textsuperscript{1–3,25} for samples up to about 200–250\,μm in thickness and a more recent\textsuperscript{6} 1-mm protocol, compared to the protocol reported here (gray boxes) enabling gradient-free en bloc staining of samples 2\,mm to 1 centimeter in size (sufficient to cover the depth of the human cortex (Fig. 2c–e) and entire adult mouse brains). \textbf{b}, Staining of a sample 2\,mm in size using the 1-mm protocol\textsuperscript{6} (top) and our protocol (bottom) (SEM images at a voxel size of 11.24\,nm; dose, 47\,e\textsuperscript{−}\,nm\textsuperscript{−2}). \textbf{c}, High-contrast staining of whole mouse brains. High-resolution EM images acquired from the superficial cortex to subcortical tissue; approximate positions are indicated in the μCT overview (left). High-resolution SEM images (at a voxel size of 5.62\,×\,5.62\,nm\textsuperscript{2}; 113\,e\textsuperscript{−}\,nm\textsuperscript{−2}). μCT data of whole brains: https://wklink.org/7742 and https://wklink.org/7789. \textbf{d}, Illustrations and examples of key challenges of hemisphere and mouse whole-brain staining. Yellow arrows point to sites of intracellular overextraction with loss of cytosolic components (left) and breakages of the sample at microscopic (middle) or macroscopic (right) scale. \textbf{e}, ATUM cutting test at the center of a fully stained and fully embedded mouse whole-brain sample (sample W3, Spurr’s resin, Supplementary Table 1). Middle, sketch and picture of ATUM cutting at sample center, trimmed to about 4\,×\,2\,mm\textsuperscript{2} for ATUM cutting. Right, sketch of the location of two artifact-free ATUM test series at the exposed surface and after an additional depth of 1\,mm (Supplementary Video 1).
\alttext{Figure 1 fills the upper two thirds of p.\,2; five panels, no plotted data.
(a) Schematic of sample sizes (dark grey blocks drawn as boxes with a bevelled top, each with a pale glow in the middle). Top row, left to right: a thin flat slab with wavy arrows labelled ``OsO\textsubscript{4} + FeCN'' and ``TCH'' entering from above and a thickness bracket ``<200 µm'', captioned ``Reduced-osmium-TCH protocol''; arrow to a taller block with a height arrow ``<\textasciitilde 1 mm'', captioned ``1-mm protocol Hua et al.\textsuperscript{6}''; arrow labelled ``×8'' to a large block with a height arrow ``2 mm'' and the grey box ``This protocol''. A diagonal arrow labelled ``×500'' points from the 1-mm block down to an outline drawing of an adult mouse brain (side view) with dimension arrows ``\textasciitilde 5–7 mm'' (height), ``\textasciitilde 10 mm'' (length) and ``\textasciitilde 10 mm'' (diagonal, at the cerebellum end), above the grey box ``Adult mouse whole brain''. Bottom row: a small block labelled ``2 mm'', an arrow to a taller block with a height arrow ``4 mm'' above the grey box ``Human cortex'', and an arrow on to the mouse brain.
(b) Two rows of SEM images of a 2-mm sample. Top row ``1-mm protocol, 2-mm sample'': an ``Overview'' block-face image (scale bar 500 µm) with three white crosses along a line from the edge to the centre, and high-resolution images ``Periphery'', ``Intermediate'' and ``Core'' (scale bar 500 nm). The periphery shows dark, sharply outlined membranes and mitochondria; the intermediate image is greyer with weaker membranes; the core image is a nearly uniform grey speckle with faint, hardly visible membranes and a dark crack along its left edge. Bottom row ``This protocol, 2-mm sample'': overview (scale bar 500 µm, three crosses) and Periphery / Intermediate / Core images (scale bar 500 nm) that all show crisp dark membranes, round mitochondria and vesicles with similar contrast from periphery to core.
(c) ``Adult mouse whole brain'': µCT views (sagittal, horizontal and coronal sections; label ``µCT''; scale bar 2 mm), bright, with a row of small white crosses in the coronal section from the cortical surface towards the interior. Five SEM images titled ``Cortex (layer 2/3)'', ``Layer 4'', ``Layer 5'', ``White matter'' and ``Striatum'' (scale bars; 500 nm printed on the last), all with dark, well-defined membranes; the white-matter image is dominated by dark myelin sheaths around elongated axons, the striatum image by many small myelinated and unmyelinated axon profiles.
(d) ``Typical artifacts'': two sketches of a brain, one with a bright centre and dark rim (``Insufficient stain penetration'') and one with a dark centre and pale rim (``Peripheral overexposure''), next to an SEM image (scale bar 1 µm) of tissue with pale, extracted-looking cytoplasm in which four yellow arrows point to extracted sites; a sketch of a tissue piece with cracks (``Microscopic breakages'') next to an SEM image of myelinated axons (scale bar 5 µm) in which a yellow arrow points to a thin white crack; a sketch of a sample broken into several pieces (``Macroscopic breakages'') next to a µCT image of a hemisphere (scale bar 2 mm) in which a yellow arrow points to a dark crack.
(e) ATUM test: ``Full-volume embedded adult mouse whole brain'' drawn as a black brain silhouette in an orange resin block with a red dashed cut line, and a curved arrow to the trimmed half block; a pencil-style sketch of an ATUM ultramicrotome with arrows labelled ``Sample'', ``Diamond knife'' and ``Ultrathin section'' on a tape; a photograph of the cutting with ``Diamond knife'' and ``Ultrathin section'' marked (scale bar 0.5 mm); at right a box-shaped sample sketch with two thin stacks of lines below it labelled ``\textasciitilde 10 µm at 37 nm'' and ``\textasciitilde 10 µm at 38 nm'' and an arrow ``1 mm trimmed'' between them.}

initial osmium tetroxide (OsO\textsubscript{4}) incubation to 3 d and 6\,d, which resulted in increased membrane contrast, also in the sample center. But, we encountered a serious obstacle not found in smaller samples: when incubating for a sufficient duration for OsO\textsubscript{4} to reach the center of the hemisphere sample, the outer parts of the sample were showing \pdfpage{3}strong signs of ultrastructural disintegration (Fig. 1d; of note, these were not apparent in µCT, only under EM imaging). The cytosol of most neurites and cell bodies had an extracted appearance, possibly indicating the removal of intracellular proteins (Fig. 1d; note that Hayat\textsuperscript{22} had already discussed this in the context of overfixation of tissue by OsO\textsubscript{4}). As the outer part of the sample corresponds to the cortical gray matter, a key target of whole-brain connectomic analysis, whole-brain staining with good stain penetration but insufficient ultrastructural \pdfpage{4}quality in the cortical periphery would be inadequate. Because the ultrastructural alterations were likely a result of protein overoxidization by prolonged OsO\textsubscript{4} incubation (Supplementary Results), which could be slowed down with a lower incubation temperature (4\,°C)\textsuperscript{22}, we used lower temperatures for OsO\textsubscript{4} and ferrocyanide (FeCN) incubation, which solved the peripheral ultrastructural disintegration problem. Additional challenges from the FeCN step were addressed by modified washing steps around the relevant incubations (Supplementary Results and Extended Data Figs. 4b–f and 5a,b).

A major remaining problem was that larger samples were consistently broken during the staining process, either in their entirety into several smaller pieces or with microbreakages that would hamper dense circuit reconstruction (Fig. 1d). These breakages would usually occur in relation to the Pg incubation step (Extended Data Fig. 4g,h). We therefore had to assess the effect of all earlier protocol steps on the stability of the sample during Pg incubation. This was of particular importance when the Pg incubation (including the washing steps before and after) was performed in water, as determined for the 2-mm samples (Extended Data Fig. 1g), and this water incubation could induce substantial osmotic forces in the sample. We found that inserting an additional extended OsO\textsubscript{4} incubation step after the FeCN step would stabilize samples such that they could be incubated in water for up to 100\,h without the occurrence of major breakages (Extended Data Fig. 4i; a possible explanation is the membrane perforation induced by OsO\textsubscript{4}, which increases tolerance to osmotic stress\textsuperscript{22}). While this additional OsO\textsubscript{4} step may also serve to enhance background staining and thus sample conductivity for SEM, in our view, the enhanced stability of large samples is the most critical aspect (see Extended Data Fig. 6h–n for an alternative approach to avoid incubation with pure water by omitting the second osmium step).

We applied the hemisphere protocol to \emph{n}\,=\,4 hemispheres (Extended Data Fig. 3a–d), which all remained intact and provided homogeneous high-contrast staining throughout the sample. In addition, we stained two entire mouse brains with modest additional extensions of incubation times (Fig. 1c and Extended Data Fig. 3e), yielding a mouse whole-brain staining protocol for large-scale connectomics\textsuperscript{11} (Extended Data Fig. 7). We observed two types of artifacts that remained in large-sample staining (Extended Data Fig. 6): the detachment of larger blood vessel walls from the surrounding neuropil and rare remaining microbreakages in the subcortical regions with a high rate of myelinated fibers. Also, special care had to be taken to preserve the integrity of the cerebellum (Extended Data Fig. 6f,g). In general, we recommend using longer protocol durations when in doubt about which protocol version to choose (for example, use the whole-brain protocol for samples even slightly larger than a hemisphere).

To confirm that our whole-brain embedding protocol would also yield sufficient sample stability for ultrathin cutting, we performed a test series of automated tape-collecting ultramicrotome (ATUM) cutting\textsuperscript{23} on the center of a mouse whole-brain sample that was stained and embedded as a whole and only trimmed for the experiment afterward (Fig. 1e and Supplementary Video 1).

To quantitatively determine the sufficiency of our staining quality and resin stability in the center of our samples for three-dimensional

\begin{center}
\includegraphics[width=0.95\linewidth,height=0.55\textheight,keepaspectratio]{figures/fig2.png}
\end{center}
\noindent\textbf{Fig. 2 | Traceability assessment in a mouse hemisphere sample and application to the human cortex.} \textbf{a}, SBEM dataset (pixel size, 11.24\,×\,11.24\,nm\textsuperscript{2}; cutting thickness, 35\,nm; electron dose, 28\,e\textsuperscript{−}\,nm\textsuperscript{−2}) from the center of a fully stained and fully embedded mouse hemisphere (H6), located in the striatum. Twenty randomly seeded axons (from a bounding box 3\,×\,3\,×\,3\,μm in size) manually reconstructed in the entire dataset. Control SBEM dataset of a mouse striatum from a conventional sample 1\,×\,1\,×\,0.5\,mm\textsuperscript{3} in size stained using the 1-mm protocol\textsuperscript{6} and analyzed in the same way (gray bar), yielding a similar mean inter-error distance for axon reconstruction (mean and 25th and 75th percentiles; Kolmogorov–Smirnov test, \emph{P}\,=\,0.4814; Methods). \textbf{b}, Automated dense axon reconstruction (method as in ref. 24) with resulting axon path-length distributions (bottom) for hemisphere striatum (blue) and control striatum (black) datasets. \textbf{c}, Immersion fixation, vibratome slicing and sample generation from human brain biopsies from neurosurgical samples (same procedure as in ref. 24). \textbf{d}, Processing of a human brain sample using this protocol for large-volume staining (here, a 2\,×\,2\,×\,3-mm block\,=\,12\,mm\textsuperscript{3}). \textbf{e}, Overview EM of a tissue block 3\,mm in height spanning the pia and layers 1–6 of the human cortex (top left); high-resolution EM images of various locations along the cortical axis (right and bottom). Pixel size in \textbf{e}, (11.24\,nm)\textsuperscript{2}; electron dose, 42\,e\textsuperscript{−}\,nm\textsuperscript{−2}.
\alttext{Figure 2 fills most of p.\,3 (its caption is at the top of p.\,4).
(a) Left: 3D rendering of the SBEM volume ``Striatum, from hemisphere'' as a box with edge lengths 73 µm (height), 50 µm and 40 µm; its faces show dense grey neuropil with dark membranes, a few large pale cell bodies with nuclei at the upper and right edges, and pale blood-vessel profiles. Top right: ``Manual axon tracing'': a wire-frame box with about 20 coloured skeleton tracings radiating in all directions from a small black seeding box in the centre, labelled ``Randomly seeded''. Bottom right: bar chart, y axis ``Mean inter-error distance along axons (mm)'', linear, 0 to 0.7 (ticks every 0.1, labels 0, 0.2, 0.4, 0.6), title ``Striatum''; blue bar ``Hemi-sphere'' about 0.47 mm with an error bar from about 0.28 to 0.67; grey bar ``Control'' about 0.49 mm with an error bar from about 0.29 to 0.68.
(b) Top: ``Automated axon reconstruction'': a box densely filled with thousands of multi-coloured reconstructed axon fragments, labelled ``Dense''. Bottom: histogram drawn as stair lines, y axis ``Automatically reconstructed axons'', logarithmic (labels 10\textsuperscript{1}, 10\textsuperscript{3}, 10\textsuperscript{5}), x axis ``Axon path length (µm)'' 0 to 100; ``Hemisphere'' (blue) and ``Control'' (black, arrow) both start above 10\textsuperscript{5} in the first bin at 0, fall steeply to about 10\textsuperscript{4} by a few µm and then decline roughly smoothly to about 10 at 70--100 µm; the two curves overlap up to about 50 µm, beyond which the control (black) lies somewhat lower and is noisier.
(c) Schematic: drawing of a human brain (scale bar 5 cm) with a red box on the frontal lobe and a red arrow to ``Biopsy'' (a block, scale bar 5 mm), arrow to ``Vibratome sections'' (two slices), arrow to ``Subsampling'' (a slice with four red rectangles).
(d) Schematic: a block with edges ``2 mm'', ``2 mm'', ``3 mm'', arrow ``2–3-mm protocol'' to a grey (stained) block with a red dashed line, arrow ``Cut in half'' to two half blocks; an eye labelled ``SEM (e)'' looks at the cut face of one half, and a curved arrow leads to a trimmed piece labelled ``ATUM''.
(e) Left: overview EM of the human tissue block (black background, scale bar 1 mm), labelled ``Pia'' at the top and ``White matter'' at the bottom, with the positions 1, 2, 3, 4, 5 marked from top to bottom; the block is evenly grey from top to bottom, with a short bright streak near position 5. Right and bottom: five high-resolution SEM images titled 1 to 5, each with a scale bar (1 µm printed on image 5), all showing dense neuropil with crisp dark membranes, mitochondria and myelinated axons, with similar contrast at all positions.}
\figuredata{a}{figures/fig2a.csv}{figures/fig2a.py}{source data (tier 1): Edmond 10.17617/3.RG58DU, data/Fig.2/Fig.2a/data.mat (per-axon path lengths and discontinuities, saved means and resampled inter-error distances); checked against the vector paths of the printed bars (bar heights agree within 0.002 mm; the printed error bars come from a different random resampling than the saved one, see CSV header)}
\figuredata{b (bottom)}{figures/fig2b.csv}{figures/fig2b.py}{source data (tier 1): Edmond 10.17617/3.RG58DU, data/Fig.2/Fig.2b/Fig2b\_data.mat (path length of every automatically reconstructed axon, H6 and control striatum), binned like MATLAB's automatic histogram; checked against the vector stair outlines of the printed figure (agree to the line width)}

(3D) EM acquisition and axon reconstruction, we acquired a serial block-face SEM (SBEM) dataset from the center (striatum) of one stained and full-volume embedded mouse brain hemisphere (Extended Data Fig. 3c; H6, webKnossos link) and compared this to a control SBEM dataset from the striatum obtained from a sample 1\,×\,1\,×\,0.5\,mm\textsuperscript{3} in size stained with the 1 mm protocol\textsuperscript{6}. Both manual axon reconstruction (Fig. 2a) and automated axon reconstruction (Fig. 2b) yielded similar results between locally stained control and hemisphere-stained samples.

Finally, we applied our protocol to a sample of human brain biopsy to directly test the applicability of our protocol to large human brain samples (Fig. 2c–e; tissue was obtained as described in ref. 24, patient H6). We stained and embedded a sample 2\,×\,2\,×\,3\,mm\textsuperscript{3} in size (that is, 12\,mm\textsuperscript{3} in volume) spanning all cortical layers from the human frontal cortex. The tissue was homogeneously stained across layers (Fig. 2e). Furthermore, we obtained a cutting series using ATUM to test cuttability for large-scale 3D EM (2,733 slices at a cutting thickness of 35–40\,nm from the sample center).

We envision that the protocols reported here will be useful for large-scale connectomic projects in mice and other species. In particular, the volume of a mouse brain also approaches relevant cortical and subcortical volumes in higher mammals such as non-human primates and humans. The goal of obtaining large connectomes from fractions or the entirety of the human cortex\textsuperscript{12} will also profit from the advances described here, both for fundamental research and possible clinical applications, such as connectomic analysis of pathological brain specimens from neurosurgical interventions\textsuperscript{24}.

\subsection*{Online content}
Any methods, additional references, Nature Portfolio reporting summaries, source data, extended data, supplementary information, acknowledgements, peer review information; details of author contributions and competing interests; and statements of data and code availability are available at https://doi.org/10.1038/s41592-023-01866-3.

\subsection*{References}
\noindent 1.\ Karnovsky, M. J. Use of ferrocyanide-reduced osmium tetroxide in electron microscopy. In \emph{11th Annual Meeting, the American Society for Cell Biology} (ASCB, 1971; 146).\par
\noindent 2.\ de Bruijn, W. C. Glycogen, its chemistry and morphologic appearance in the electron microscope I. A modified OsO\textsubscript{4} fixative which selectively contrast glycogen. \emph{J. Ultrastruct. Res.} \textbf{42}, 29–50 (1973).\localcopy{debruijn1973}\par
\noindent 3.\ Willingham, C. \& Rutherford, A. Use of osmium-thiocarbohydrazide-osmium and ferrocyanide-reduced methods to enhance membrane contrast and preservation in cultured cells. \emph{J. Histochem. Cytochem.} \textbf{32}, 455–460 (1984).\localcopy{willingham1984}\par
\noindent 4.\ Briggman, K. L., Helmstaedter, M. \& Denk, W. Wiring specificity in the direction-selectivity circuit of the retina. \emph{Nature} \textbf{471}, 183–190 (2011).\par
\noindent 5.\ Holcomb, P. S. et al. Synaptic inputs compete during rapid formation of the calyx of Held: a new model system for neural development. \emph{J. Neurosci.} \textbf{33}, 12954–12969 (2013).\par
\noindent 6.\ \pdfpage{5} Hua, Y., Laserstein, P. \& Helmstaedter, M. Large-volume en-bloc staining for electron microscopy-based connectomics. \emph{Nat. Commun.} \textbf{6}, 1–7 (2015).\localcopy{hua2015}\par
\noindent 7.\ Motta, A. et al. Dense connectomic reconstruction in layer 4 of the somatosensory cortex. \emph{Science} \textbf{366}, eaay3134 (2019).\par
\noindent 8.\ Hayworth, K. J. et al. Gas cluster ion beam SEM for imaging of large tissue samples with 10\,nm isotropic resolution. \emph{Nat. Methods} \textbf{17}, 68–71 (2020).\par
\noindent 9.\ Kuan, A. T. et al. Dense neuronal reconstruction through X-ray holographic nano-tomography. \emph{Nat. Neurosci.} \textbf{23}, 1637–1643 (2020).\par
\noindent 10.\ Yin, W. et al. A petascale automated imaging pipeline for mapping neuronal circuits with high-throughput transmission electron microscopy. \emph{Nat. Commun.} \textbf{11}, 4949 (2020).\par
\noindent 11.\ Abbott, L. F. et al. The mind of a mouse. \emph{Cell} \textbf{182}, 1372–1376 (2020).\par
\noindent 12.\ Motta, A. et al. Big data in nanoscale connectomics, and the greed for training labels. \emph{Curr. Opin. Neurobiol.} \textbf{55}, 180–187 (2019).\par
\noindent 13.\ Mikula, S. \& Denk, W. High-resolution whole-brain staining for electron microscopic circuit reconstruction. \emph{Nat. Methods} \textbf{12}, 541–546 (2015).\localcopy{mikula2015}\par
\noindent 14.\ Januszewski, M. et al. High-precision automated reconstruction of neurons with flood-filling networks. \emph{Nat. Methods} \textbf{15}, 605–610 (2018).\par
\noindent 15.\ Song, K., Feng, Z. \& Helmstaedter, M. High-contrast en-bloc staining of mouse whole-brain samples for EM-based connectomics. Preprint at \emph{bioRxiv} 10.1101/2022.03.30.486341 (2022).\par
\noindent 16.\ Ströh, S. et al. In situ X-ray-assisted electron microscopy staining for large biological samples. \emph{eLife} \textbf{11}, e72147 (2022).\localcopy{stroh2022}\par
\noindent 17.\ Kazmierski, T. M. \& Millington, P. F. Modification of an SEM for backscattered electron imaging at low vacuum. \emph{J. Microsc.} \textbf{117}, 431–436 (1979).\par
\noindent 18.\ Sommons, R. \& Marquis, P. Application of the low vacuum scanning electron microscope to the study of biomaterials and mammalian cells. \emph{Biomaterials} \textbf{18}, 81–86 (1997).\par
\noindent 19.\ Vidil, T. et al. Control of reactions and network structures of epoxy thermosets. \emph{Prog. Polym. Sci.} \textbf{62}, 126–179 (2016).\par
\noindent 20.\ Korogod, N., Petersen, C. C. \& Knott, G. W. Ultrastructural analysis of adult mouse neocortex comparing aldehyde perfusion with cryo fixation. \emph{eLife} \textbf{4}, e05793 (2015).\par
\noindent 21.\ Loos, M. R. et al. The effect of acetone addition on the properties of epoxy. \emph{Polimeros} \textbf{18}, 76–80 (2008).\par
\noindent 22.\ Hayat, M. A. \emph{Fixation for Electron Microscopy} (Academic Press, 1981).\par
\noindent 23.\ Hayworth, K. J. et al. Automating the collection of ultrathin serial sections for large volume TEM reconstructions. \emph{Microsc. Microanal.} \textbf{12}, 86–87 (2006).\par
\noindent 24.\ Loomba, S. et al. Connectomic comparison of mouse and human cortex. \emph{Science} \textbf{377}, eabo0924 (2022).\par
\noindent 25.\ Seligman, A. M., Wasserkrug, H. L. \& Hanker, J. S. A new staining method (OTO) for enhancing contrast of lipid-containing membranes and droplets in osmium tetroxide-fixed tissue with osmiophilic thiocarbohydrazide (TCH). \emph{J. Cell Biol.} \textbf{30}, 424–432 (1966).\localcopy{seligman1966}\par

\noindent\textbf{Publisher’s note} Springer Nature remains neutral with regard to jurisdictional claims in published maps and institutional affiliations.

\noindent\textbf{Open Access} This article is licensed under a Creative Commons Attribution 4.0 International License, which permits use, sharing, adaptation, distribution and reproduction in any medium or format, as long as you give appropriate credit to the original author(s) and the source, provide a link to the Creative Commons license, and indicate if changes were made. The images or other third party material in this article are included in the article’s Creative Commons license, unless indicated otherwise in a credit line to the material. If material is not included in the article’s Creative Commons license and your intended use is not permitted by statutory regulation or exceeds the permitted use, you will need to obtain permission directly from the copyright holder. To view a copy of this license, visit http://creativecommons.org/licenses/by/4.0/.

\noindent © The Author(s) 2023 \pdfpage{6}

\section*{Methods}

\subsection*{Animal experiments}
All experimental procedures were approved by the local animal care and use committee and were in accordance with the laws of animal experimentation issued by the German federal government (Regierungspräsidium Darmstadt, Germany, permits V54, 19 c 20/15, F126/1028 and F126/1002).

Adult C57BL6/J mice (male and female, P30–P90) were treated with analgesics (0.1\,mg per kg buprenorphin (CP-Pharma) and 100\,mg per kg metamizol (WDT)) for 0.5\,h before isoflurane anesthesia (Harvard Apparatus; 5\% in O\textsubscript{2} for initialization, 2–3\% for maintenance; O\textsubscript{2} flow rate, 1\,l\,min\textsuperscript{−1}). Following anesthesia, the animals were transcardially perfused (Harvard Apparatus, flow rate of 10\,ml\,min\textsuperscript{−1}) using 15\,ml sodium cacodylate buffer (0.15\,M, pH 7.4, Sigma-Aldrich) followed by 30\,ml fixative containing 2.5\% paraformaldehyde (Sigma-Aldrich), 1.25\% glutaraldehyde (Serva) and 2\,mM calcium chloride (Sigma-Aldrich) in 0.08\,M sodium cacodylate buffer (osmolarity of about 700–800\,mmol per kg, pH 7.4). The duration between the start of perfusion and incision of the diaphragm was less than 30\,s. After perfusion, animals were decapitated, the skull was opened with care to avoid mechanical damage to the brain, and the brain was post-fixed in situ for 12–96\,h at 4\,°C before extraction from the skull.

For 2-mm samples, the brain was cut into 2-mm-thick coronal sections in 0.15\,M cacodylate buffer using a vibratome (Leica, VT1200). Next, a biopsy punch 2\,mm in diameter (Kai Medical) was used to extract samples from dorsal cortical, ventral cortical and subcortical regions (Supplementary Fig. 1h). The samples were then stored in 0.15\,M cacodylate buffer at 4\,°C for 8–24\,h before staining.
\decompnote{There is no Supplementary Fig.\ 1h: the Supplementary Information has no figures. The sampling positions are shown in Extended Data Fig.\ 1h (photograph with positions 1--3).}

For hemisphere samples, brains were cut with a razor blade (Wilkinson) along the midline. The hemispheres were stored in 0.15\,M cacodylate buffer at 4\,°C for 24\,h before staining.

\subsection*{Human sample}
The human brain tissue sample (human individual H6, same as in ref. 24) was collected during a neurosurgical procedure that was indicated for medical reasons and independently from this research project at the Department of Neurosurgery at the Klinikum rechts der Isar of the Technical University of Munich. The sample was obtained from access tissue (presumably healthy brain parenchyma that had to be removed as part of the procedure and would have been discarded otherwise) before removal of the respective target lesion, as approved by the Ethics Committee of the Technical University of Munich School of Medicine (Ethikvotum 184/16S and 273/21 S-EB). All patients had given their written informed consent.

The human ‘H6’ sample was obtained from the inferior frontal gyrus from a 69-year old female patient during surgical removal of a frontal mass lesion (final diagnosis, glioblastoma multiforme), the same patient as in ref. 24. Following surgical removal, tissue was directly collected in fix solution kept at 4\,°C. The tissue was immediately sliced into 2-mm-thin slices in cold fixative using a vibratome. Slices were kept at 4\,°C overnight. Samples sized 2\,×\,3\,mm\textsuperscript{2} were then cut out from the slices with razor blades.

The 2\,×\,3\,×\,2-mm\textsuperscript{3} samples were stained based on the 2-mm protocol (Extended Data Fig. 7) with minor modifications (Supplementary Table 1). In brief, after four rinses for 30\,min with 0.15\,M sodium cacodylate buffer (CaC) (4\,°C), the following steps were applied sequentially: 22\,h, 2\% OsO\textsubscript{4} in 0.15\,M CaC (pH 7.4) at room temperature (RT); four rinses for 30\,min with 0.15\,M CaC at RT; 22\,h, 2.5\% FeCN in 0.15\,M CaC (pH 7.4) at RT; 6\,h, 2\% OsO\textsubscript{4} in 0.15\,M CaC (pH 7.4) at RT; 30\,min, 0.15\,M CaC; three rinses for 20\,min with water at RT; 18\,h, 2\% Pg in water at RT; three rinses for 20\,min with water; 6\,h, 2\% OsO\textsubscript{4} in water; three rinses for 20\,min with water; 14\,h, 4\% uranyl acetate (UA) at 4\,°C (afterward, switch temperature to 50\,°C for 2\,h); three times for 20\,min with water. Afterward, the samples went through graded ethanol dehydration steps: 30\,min, 20\% ethanol at 4\,°C; 30\,min, 40\% ethanol at 4\,°C; 30\,min, 60\% ethanol at 4\,°C; 30\,min, 80\% ethanol at 4\,°C; 45\,min, 100\% ethanol at RT. For the resin infiltration, after three rinses for 45\,min with pure acetone, the samples went through graded Epon resin steps (for 10\,ml resin: 5.9\,g Epon medium, 2.25\,g DDSA, 3.7\,g MNA, 205\,µl DMP) in acetone, all at 4\,°C: 4\,h, 12.5\%; 13\,h, 25\% (overnight); 4\,h, 37.5\%; 4\,h, 50\%; 19\,h, 62.5\% (overnight); 8\,h, 75\%; 19\,h, 87.5\% (overnight); 8\,h, 95\%; 19\,h, 95\% (first overnight for 95\%); 8\,h, 95\%; 19\,h, 95\% (second overnight for 95\%); 8\,h, 95\%; 19\,h, 95\% (third overnight for 95\%); 8\,h, 100\%; 19\,h, 100\% (first overnight for 100\%); 8\,h, 100\%; 19\,h, 100\% (second overnight for 100\%); afterward, the samples were embedded in freshly prepared Epon resin and cured at 60\,°C for 3\,d.

\subsection*{Staining experiments}
All staining and resin-infiltration steps for 2-mm samples were carried out in 2-ml Eppendorf tubes and, for hemisphere or whole-brain samples, in 50-ml glass tubes at RT (\textasciitilde{}20–23\,°C) unless specified otherwise. For steps involving photosensitive chemicals (FeCN, TCH, Pg, UA or lead aspartate (Ld)), the tubes were covered with aluminum foil.

All chemicals used in the staining pipeline are listed in Supplementary Table 2; all experiments reported in the main and supplementary figures are detailed in Supplementary Table 3.

For simplicity, the following terms were used to refer to the recurring staining steps in different experiments: CaC, sodium cacodylate buffer rinse, 0.15\,M, pH 7.4; water, water (Milli-Q) rinse; first Os, the first OsO\textsubscript{4} incubation step, 2\% OsO\textsubscript{4} in 0.15\,M CaC, pH 7.4; FeCN, the FeCN incubation step, 2.5\% in 0.15\,M CaC, pH 7.4; second Os, the OsO\textsubscript{4} incubation step after FeCN, 2\% OsO\textsubscript{4} in 0.15\,M CaC, pH 7.4; Pg, the Pg incubation step, 4\% in water; TCH, the TCH incubation step, 1\% in water; third Os, the OsO\textsubscript{4} incubation step after Pg, 2\% in water; UA, the UA incubation step, UA (2\% or 4\%) dissolved in water; Ld, the Ld incubation step, 0.66\% lead nitrate in 0.03\,M aspartic acid, pH 5.0.

\subsection*{Staining of 2-mm samples with the 1-mm protocol}
The 1-mm staining protocol (ref. 6 with addition of the second Os step) was applied to 2-mm samples directly. The samples were stained with the following steps: first Os, 1.5\,h; FeCN, 1.5\,h; second Os, 1\,h; CaC, 0.5\,h; water, 0.5\,h; TCH, 1.5\,h; water, 0.5\,h (two times); third Os, 1.5\,h; water, 0.5\,h (two times); 2\% UA (17\,h at 4\,°C, 50\,h at 50\,°C); water, 0.5\,h (two times); Ld, 2\,h at 50\,°C; water, 0.5\,h (two times). Afterward, the samples were incubated in a graded ethanol series from 50\% (4\,°C) to 75\% (4\,°C) to 100\% (each step, 45\,min). Next, the samples were incubated in pure acetone three times, each time for 45\,min. Afterward, they were incubated in 50\% Spurr’s resin (Sigma-Aldrich, at a ratio of 0.95\,g ERL-4221, 5.9\,g DER 736, 0.1\,g NSA, 113\,µl DMAE) in acetone for 6\,h. Next, the samples were left overnight with the cap of the Eppendorf tubes open to allow evaporation of acetone. Afterward, the samples were transferred into pure Spurr’s resin for 6\,h before embedding and curing at 70\,°C for 1–3\,d. After resin curing, the samples were first imaged by μCT for staining homogeneity. Next, they were trimmed to expose the sample center; the sample surface was smoothed and imaged in high-vacuum SEM and energy-dispersive X-ray spectroscopy (EDS).
\decompnote{Two numbers in this paragraph look like typos (kept as printed): ``50\,h at 50\,°C'' for UA (Supplementary Table 3, experiment 1, and the other Methods sections give 2\,h at 50\,°C); and the Spurr's resin recipe ``0.95\,g ERL-4221, 5.9\,g DER 736, 0.1\,g NSA, 113\,µl DMAE'', repeated in the Spurr's sections below, differs from Supplementary Table 1 (4.1\,g ERL 4221, 0.95\,g DER 736, 5.9\,g NSA, 113\,μl DMAE).}

\subsection*{Step-by-step μCT diagnosis of the main staining steps of the 1-mm protocol on 2-mm samples}
The 1-mm staining protocol (ref. 6 with the addition of the second Os step) was applied to a batch of 2-mm samples as described above: first Os, 1.5\,h; FeCN, 1.5\,h; second Os, 1\,h; CaC, 0.5\,h; water, 0.5\,h; TCH, 1.5\,h; water, 0.5\,h (two times); third Os, 1.5\,h; water, 0.5\,h (two times); 2\% UA (17\,h at 4\,°C, 50\,h at 50\,°C); water, 0.5\,h (two times); Ld, 2\,h at 50\,°C; water, 0.5\,h (two times). After each main step (first Os, FeCN, second Os, third Os, UA), two samples were taken out from the staining pipeline and rinsed with either CaC or water, depending on the solvent condition of the corresponding staining step (rinsing solution was changed every 8 or 17\,h). When the rinsing was complete for all samples from different conditions, samples were embedded in 2\% agarose in water \pdfpage{7}in the same Eppendorf tube at different tube depths and stored at 4\,°C for the agarose to cure. Next, the tube containing all the samples was imaged by μCT to investigate staining gradients.

\subsection*{Extending the FeCN incubation time for 2-mm samples}
The 2-mm samples were stained as follows: first Os, 3\,h; FeCN. During the FeCN incubation, two samples were taken out at each of the following time points: 1.5\,h, 3\,h, 7\,h and 17\,h. Samples were then rinsed with CaC for 1\,h and embedded in 2\% agarose as described above. After agarose curing, they were imaged by μCT to investigate staining gradients possibly related to FeCN incubation.

\subsection*{Interaction between OsO\textsubscript{4} incubation duration and FeCN incubation}
The 2-mm samples were stained using the following two conditions: (1) first Os, 3\,h; FeCN, 17\,h or (2) first Os, 24\,h; FeCN, 17\,h. Next, they were rinsed sequentially with CaC (for 0.5\,h) and water (two times, each time for 0.5\,h); they were then dehydrated and embedded according to the 2-mm Spurr’s resin protocol. After resin embedding, they were trimmed to expose the center, their surface was smoothed, and they were imaged by low-vacuum SEM to investigate membrane contrast.

\subsection*{Extending the TCH incubation for 2-mm samples}
Three groups of 2-mm samples were stained; each batch was incubated with TCH for a different incubation length, using the following protocol: first Os, 3\,h; FeCN, 17\,h; second Os, 3\,h; CaC, 0.5\,h; water, 0.5\,h (two times); TCH for 1.5\,h or 3\,h or 5\,h; water, 0.5\,h (two times); third Os, 3\,h. Afterward, they were rinsed and embedded in 2\% agarose for imaging by μCT.

\subsection*{Replacing TCH with pyrogallol}
Three groups of 1-mm samples were stained; each group differed in the TCH-related step, using the following protocol: first Os, 1.5\,h; FeCN, 1.5\,h; second Os, 1\,h; CaC, 0.5\,h; water, 0.5\,h; TCH, 1.5\,h or Pg, 1.5\,h or water, 1.5\,h; water, 0.5\,h (two times); third Os, 1.5\,h; water, 0.5\,h (two times); 2\% UA (17\,h at 4\,°C, 50\,h at 50\,°C); water, 0.5\,h (two times); Ld, 2\,h at 50\,°C; water, 0.5\,h (two times). Afterward, they were dehydrated and embedded in resin according to ref. 6. After resin curing, the samples were trimmed to expose the center, their surface was smoothed, and they were imaged by high-vacuum SEM and EDS. For EDS measurement, we selected point measurements in the neuropile of the sample center.

\subsection*{Comparison of pyrogallol incubation in water versus CaC}
Two groups of 2-mm samples (the groups differ in Pg incubation and the water steps around Pg) were stained with the following sequential incubation steps: first Os, 24\,h; FeCN, 17\,h; second Os, 3\,h; CaC, 0.5\,h; water, 0.5\,h (twice); Pg in water, 17\,h; water, 0.5\,h (twice) or CaC, 0.5\,h (twice); Pg in CaC, 17\,h; CaC, 0.5\,h (twice); third Os, 6\,h; water, 0.5\,h (twice); 4\% UA (17\,h at 4\,°C, 2\,h at 50\,°C); water, 0.5\,h (twice). Samples were then dehydrated and embedded in Spurr’s resin as described above for 2-mm samples. After resin curing, they were trimmed to expose the center, their surface was smoothed, and they were imaged by high-vacuum SEM.

\subsection*{Long-duration OsO\textsubscript{4} incubation for 2-mm samples}
The 2-mm samples were stained with 2\% OsO\textsubscript{4} in CaC for either 3\,d or 6\,d. After rinsing with 0.15\,M CaC for 0.5\,h and two times with water (each time for 0.5\,h), the samples were dehydrated and embedded according to the 2-mm Spurr’s resin protocol. Next, they were trimmed to expose the center, their surface was smoothed, and they were imaged by low-vacuum SEM for determining ultrastructural preservation.

\subsection*{Extended pyrogallol and the third OsO\textsubscript{4} incubation for 2-mm samples}
Three groups of 2-mm samples (groups differ in the incubation duration of Pg and third Os steps) were stained using the following protocol: first Os, 3\,h; FeCN, 17\,h; second Os, 3\,h; CaC, 0.5\,h; water, 0.5\,h (twice); Pg in water (for 6 or 17\,h); water, 0.5\,h (twice); third Os (for 3 or 6\,h); water, 0.5\,h (twice); 2\% UA (17\,h at 4\,°C, 2\,h at 50\,°C); water, 0.5\,h (twice). The combinations of Pg and third Os incubation times were as follows: (1) Pg, 6\,h; third Os, 3\,h; (2) Pg, 17\,h; third Os, 3\,h; (3) Pg, 17\,h; third Os, 6\,h. After staining, the samples were dehydrated and embedded in Spurr’s resin as described above for 2-mm samples. After resin curing, they were trimmed to expose the center, their surface was smoothed, and they were imaged by high-vacuum SEM.

\subsection*{UA and Ld steps for 2-mm samples}
Four groups of 2-mm samples (groups differ in incubation of UA and Ld steps) were stained using the following protocol: first Os, 3\,h; FeCN, 17\,h; second Os, 3\,h; CaC, 0.5\,h; water, 0.5\,h (twice); Pg in water (17\,h); water, 0.5\,h (twice); third Os, 6\,h; water, 0.5\,h (twice); 2\% or 4\% UA (17\,h at 4\,°C, 2\,h at 50\,°C); water, 0.5\,h (twice); Ld, 50\,°C, 4\,h or 24\,h; water, 0.5\,h (twice). The combinations of UA and Ld incubations were as follows: (1) 2\% UA; no Ld; (2) 2\% UA; Ld, 4\,h; (3) 4\% UA; no Ld; (4) 4\% UA; Ld, 24\,h. After staining, the samples were dehydrated and embedded in Spurr’s resin as described above for 2-mm samples. After resin curing, samples were imaged by μCT and then trimmed to expose the center; their surface was smoothed and imaged by high-vacuum SEM.

\subsection*{Temperature of the OsO\textsubscript{4} incubation step}
Two groups of 2-mm samples were stained according to the following steps: (1) Os (4\,°C, 7\,d) or (2) Os (4\,°C, 7\,d; RT, 1\,d). After staining, the samples were dehydrated and embedded in Spurr’s resin as described above for 2-mm samples. After resin curing, they were trimmed to expose the center, their surface was smoothed, and they were then imaged by low-vacuum SEM.

\subsection*{OsO\textsubscript{4} incubation at 4\,°C followed by FeCN incubation}
Four groups of samples were stained according to the following steps: (1) Os (4\,°C, 6\,d; RT, 1\,d), FeCN (RT, 1\,d); (2) Os (4\,°C, 6\,d; RT, 1\,d), FeCN (4\,°C, 1\,d); (3) Os (4\,°C, 7\,d), FeCN (RT, 1\,d); (4) Os (4\,°C, 7\,d), FeCN (4\,°C, 1\,d). After staining, the samples were dehydrated and embedded in Spurr’s resin as described above for 2-mm samples. After resin curing, they were trimmed to expose the center, their surface was smoothed, and they were imaged by low-vacuum SEM.

\subsection*{Diffusion of OsO\textsubscript{4} at 4\,°C in hemisphere samples}
A hemisphere was incubated in OsO\textsubscript{4} at 4\,°C. At different time points (17\,h, 24\,h, 40\,h), it was taken out of the fridge to perform a fast μCT scan (usually about 15–20\,min in total) and then returned to the fridge at 4\,°C. µCT images were analyzed using Zeiss TXM3DViewer software, in which the depth of OsO\textsubscript{4} diffusion was measured in sagittal reslices.

\subsection*{Effect of CaC steps on the staining gradient in hemispheres}
Two groups of hemispheres were stained with the following two conditions: (1) Os, 48\,h; FeCN, 48\,h or (2) Os, 48\,h; CaC, 48\,h; FeCN, 48\,h. After staining, the hemispheres were briefly rinsed with CaC and cut into coronal sections of about 2\,mm in thickness with a razor blade. Afterward, the coronal sections were dehydrated and embedded in Spurr’s resin as described above for 2-mm samples. After resin curing, they were imaged by μCT and then trimmed flat to expose the top surface and smoothed and imaged by low-vacuum SEM.

\subsection*{Effect of CaC steps on the velocity of FeCN diffusion}
Two groups of 2-mm samples were stained with the following two conditions: (1) Os, 24\,h; FeCN, 1.5\,h or (2) Os, 24\,h; CaC, 24\,h; FeCN, 1.5\,h. After staining, the samples were dehydrated and embedded in Spurr’s resin as described above for 2-mm samples. After resin curing, they were imaged by μCT and then trimmed flat to expose the top \pdfpage{8}surface and smoothed and imaged by low-vacuum SEM and EDS. For EDS measurement, we selected point measurements in the neuropile of the sample center.

\subsection*{Interaction of pyrogallol incubation with the Os–FeCN gradient (H3)}
A hemisphere sample (H3) was stained according to the following steps: first Os (4\,°C, 96\,h; RT, 24\,h), CaC (4\,°C, 48\,h); FeCN (4\,°C, 48\,h); second Os (72\,h); CaC (24\,h); water (29\,h); Pg (24\,h); water (48\,h); third Os (48\,h); 4\% UA (4\,°C, 48\,h; 50\,°C, 5\,h); water (42\,h). For all CaC and water steps, the corresponding solutions were changed every 4\,h or overnight (that is, once in the morning, noon and afternoon). During staining, the hemispheres were imaged by μCT at the main staining steps (after FeCN, after second Os, after Pg). After staining, the hemisphere was cut into coronal sections about 2\,mm thick with a razor blade. Afterward, the coronal sections were dehydrated and embedded in Spurr’s resin as described above for 2-mm samples. After resin curing, they were trimmed flat to expose the top surface, smoothed and imaged by high-vacuum SEM.

\subsection*{Interaction of pyrogallol incubation with the Os–FeCN gradient and CaC (H1 and H2)}
Two hemisphere samples (H1 and H2) were stained according to the following steps: first Os (4\,°C, 63\,h; RT, 24\,h); CaC (4\,°C, 24\,h); FeCN (4\,°C, 48\,h); CaC (24\,h); second Os (32\,h); CaC (17\,h); water (9\,h); Pg (48\,h); water (40\,h); third Os (48\,h); 4\% UA (4\,°C, 48\,h; 50\,°C, 5\,h); water (48\,h). For all CaC and water steps, the corresponding solutions were changed every 4\,h or overnight (that is, once in the morning, noon and afternoon). During staining, the hemispheres were imaged by μCT at the main staining steps (after FeCN, after second Os, after Pg). After staining, the hemisphere was cut into coronal sections about 2\,mm thick with a razor blade. Afterward, the coronal sections were dehydrated and embedded in Spurr’s resin as described above for 2-mm samples. After resin curing, they were trimmed flat to expose the top surface, smoothed and imaged by high-vacuum SEM.

\subsection*{Effect of CaC incubation at room temperature for 2\,d on the Os–FeCN gradient (H4–H6)}
Three hemisphere samples (H4–H6) were stained according to the following steps: first Os (4\,°C, 72\,h; RT, 24\,h); CaC (48\,h); FeCN (4\,°C, 48\,h); CaC (24\,h); second Os (48\,h); CaC (24\,h); water (24\,h); Pg (24\,h); water (24\,h); third Os (24\,h); 4\% UA (4\,°C, 48\,h; 50\,°C, 5\,h); water (48\,h). For all CaC and water steps, the corresponding solutions were changed every 4\,h or overnight (that is, once in the morning, noon and afternoon). During staining, the hemispheres were imaged by μCT at the main staining steps (after FeCN, after second Os, after Pg). After staining, H5 was cut into coronal sections about 2\,mm thick with a razor blade. Afterward, the coronal sections were dehydrated and embedded in Spurr’s resin as described above for 2-mm samples. H4 and H6 were embedded in Spurr’s resin. After resin curing, H5 was trimmed flat to expose the top surface, smoothed and imaged by high-vacuum SEM. H4 and H6 were trimmed to expose the center, and the surface was smoothed and imaged by high-vacuum SEM.

\subsection*{Effect of CaC incubation at 4\,°C for 4\,d on the Os–FeCN gradient (H13)}
A hemisphere sample (H13) was stained according to the following steps: first Os (4\,°C, 68\,h); CaC (4\,°C, 96\,h); FeCN (4\,°C, 72\,h); CaC (48\,h); second Os (48\,h); CaC (24\,h); water (24\,h); Pg (24\,h); water (24\,h); third Os (41\,h); 4\% UA (4\,°C, 48\,h; 50\,°C, 12\,h); water (24\,h). For all CaC and water steps, the corresponding solutions were changed every 4\,h or overnight (that is, once in the morning, noon and afternoon). During staining, the hemispheres were imaged by μCT at the main staining steps (after the second Os step and after Pg). After staining, the hemisphere was embedded according to Spurr’s resin embedding for hemispheres. After resin curing, it was trimmed to expose the center, and the surface was smoothed and imaged by high-vacuum SEM.

\subsection*{Whole-brain staining (W1)}
A mouse whole-brain sample (W1) was stained according to the following steps: first Os (4\,°C, 96\,h); CaC (4\,°C, 168\,h); FeCN (4\,°C, 72\,h); CaC (4\,°C, 48\,h; RT, 48\,h); second Os (48\,h); CaC (72\,h); water (48\,h); Pg (24\,h); water (48\,h); third Os (96\,h); 4\% UA (4\,°C, 48\,h; 50\,°C, 5.5\,h); water (24\,h). For all CaC and water steps, the corresponding solutions were changed every 4\,h or overnight (that is, once in the morning, noon and afternoon). During staining, the brain was imaged by μCT at the main staining steps (during the first Os step, after the second Os step, after Pg). After staining, the brain was cut into coronal sections \textasciitilde{}2\,mm thick. Next, the sections were embedded according to Spurr’s resin embedding for 2-mm samples. After resin curing, the sections were trimmed flat to expose the surface, smoothed and checked by high-vacuum SEM.

\subsection*{Effect of the second Os step on sample stability in water}
Two groups of 2-mm samples (the groups differed in whether they were exposed to a second Os step) were incubated under the following conditions: (1) first Os, 24\,h; CaC, 24\,h; FeCN, 24\,h; second Os, 3\,h; CaC, 0.5\,h; water (two times, 0.5\,h); water or (2) first Os, 24\,h; CaC, 24\,h; FeCN, 24\,h; CaC, 0.5\,h; water (two times, 0.5\,h); water. Afterward, the samples were kept in water and imaged with the light microscope at different time points (from 2\,h to 100\,h) to investigate macroscopic sample integrity.

\subsection*{Pyrogallol–osmium(VI) interaction}
For Extended Data Fig. 6, mouse hemisphere samples were stained according to one of the following options: (1) Os, 4\,°C, 72\,h; CaC, 4\,°C, 96\,h; FeCN, 4\,°C, 72\,h; CaC, 48\,h; Os, 48\,h; CaC, 24\,h; CaC, 24\,h; Pg, 24\,h; (2) Os, 4\,°C, 72\,h; CaC, 4\,°C, 96\,h; FeCN, 4\,°C, 72\,h; CaC, 48\,h; CaC, 48\,h; CaC, 24\,h; CaC, 24\,h; Pg, 24\,h; (3) Os, 4\,°C, 72\,h; CaC, 4\,°C, 96\,h; FeCN, 4\,°C, 72\,h; CaC, 48\,h; Os, 48\,h; KCl, 24\,h; KCl, 24\,h; Pg, 24\,h; (4) Os, 4\,°C, 72\,h; CaC, 4\,°C, 96\,h; FeCN, 4\,°C, 72\,h; CaC, 48\,h; Os, 48\,h; KCl, 24\,h, pH 7; KCl, 24\,h, pH 7; KCl, pH 1; KCl, 24\,h, pH 7; Pg, 24\,h. After Pg, the samples were scanned by μCT with the parameters described in the section X-ray microtomography volumetric imaging of samples without resin embedding.

Afterward, the hemisphere sample from condition (2) was cut into \textasciitilde{}2-mm coronal chunks for further processing. The coronal sections were rinsed with CaC for 24\,h and then stained with either Os in water or Os in CaC for 24\,h. Afterward, they were embedded according to the 2-mm Spurr’s resin-embedding protocol. After the resins were cured, the samples were cut into halves and checked with low-vacuum SEM for membrane contrast.

\subsection*{Dehydration of 2-mm samples}
After the last water rinsing step in the staining protocol, 2-mm samples were exposed to graded dehydration series in ethanol (50\%, 4\,°C; 75\%, 4\,°C; and 100\%, RT; each step for 45\,min) and then three rounds of acetone incubation (each 45\,min).

\subsection*{Dehydration of hemispheres}
After the last water rinsing step in the staining protocol, hemisphere samples were exposed to graded dehydration series in ethanol (25\%, 4\,°C; 50\%, 4\,°C; 75\%, 4\,°C; and 100\%, RT; each step for 8\,h or overnight) and then three rounds of acetone incubation (each time for 8\,h or overnight).

\subsection*{Dehydration of whole-brain samples}
After the last water rinsing step in the staining protocol, whole-brain samples (for example, W3) for subsequent full-volume embedding were exposed to graded dehydration series in ethanol (25\% at 4\,°C for 12\,h, 50\%, 75\% and 100\%; each step at 4\,°C and for 24\,h) and then three rounds of acetone incubation (each time for 24\,h). \pdfpage{9}

\subsection*{Infiltration of Spurr’s resin and embedding of 2-mm samples}
After dehydration, a graded incubation of Spurr’s resin (0.95\,g ERL-4221, 5.9\,g DER 736, 0.1\,g NSA, 113\,µl DMAE) was applied at 25\%, 50\% and 75\% in acetone (each step for 8\,h or overnight). Afterward, the samples were incubated for four rounds in 100\% Spurr’s resin (each time for 8\,h or overnight). Finally, they were embedded in freshly prepared Spurr’s resin and cured at 70\,°C for 1–3\,d. All steps were performed at 4\,°C. For each resin exchange, the tubes were taken out from the fridge 20–30\,min before to be warmed to RT.

\subsection*{Infiltration of Spurr’s resin and embedding of hemisphere samples}
After dehydration, a graded incubation of Spurr’s resin (0.95\,g ERL-4221, 5.9\,g DER 736, 0.1\,g NSA, 113\,µl DMAE) was applied at 25\%, 50\% and 75\% in acetone (each step for 24\,h). Afterward, the samples were incubated for 2\,d in 90\% resin, for 3\,d in 95\% resin and for 3\,d in 100\% resin (changing every 8\,h or overnight for 90\%, 95\% and 100\% resin steps). Finally, they were embedded in freshly prepared Spurr’s resin and cured at 70\,°C for 1–3\,d. All resin steps were performed at 4\,°C. For each resin exchange, the tubes were taken out from the fridge 20–30\,min before to be warmed to RT.

\subsection*{Infiltration of Spurr’s resin and full-volume embedding of whole-brain samples}
After dehydration, a graded incubation of Spurr’s resin (0.95\,g ERL-4221, 5.9\,g DER 736, 0.1\,g NSA, 113\,µl DMAE) was applied at 25\%, 50\% and 75\% in acetone (25\% and 50\% for 2\,d and 75\% for 3\,d). Afterward, whole-brain samples (for example, W3) were incubated for 2\,d in 90\% resin, for 4\,d in 95\% resin and for 6\,d in 100\% resin. The solutions with 95\% and 100\% resin were changed every 2\,d. Finally, they were embedded in freshly prepared Spurr’s resin and cured at 70\,°C for 3\,d. All resin steps were performed at 4\,°C. For each resin exchange, the tubes were taken out from the fridge 20–30\,min before to be warmed to RT.

\subsection*{Measurement of fluidity of Epon resin}
Ten milliliters of Epon resin (Sigma-Aldrich, with a ratio of 5.9\,g Epon medium, 2.25\,g DDSA, 3.7\,g MNA and 205\,µl DMP) was prepared and stored in 15-ml Falcon tubes for different conditions: pure resin at RT (\emph{n}\,=\,2, green curves in Extended Data Fig. 1k), pure resin at 4\,°C (\emph{n}\,=\,3, blue curves in Extended Data Fig. 1k) and 95\% resin with 5\% acetone at 4\,°C (\emph{n}\,=\,4, red curves in Extended Data Fig. 1k). At different time points, videos were acquired of the resin moving in the tube after being turned upside down. The distance of movement of the resin in the tube was measured using the ticks (ml) printed on the Falcon tubes. The speed of resin movement was used as a measurement of the resin’s fluidity.

\subsection*{Infiltration with Epon resin and embedding of 2-mm samples}
After dehydration, samples were incubated in pure acetone three times, each time for 45\,min. Next, a graded incubation of Epon resin (5.9\,g Epon medium, 2.25\,g DDSA, 3.7\,g MNA, 205\,µl DMP) was applied at 12.5\%, 25\%, 37.5\%, 50\%, 62.5\%,75\% and 87.5\% in acetone (each step for 4\,h or overnight). Afterward, the samples were incubated for four to six rounds in 95\% Epon resin (each time for 8\,h or overnight) and then for four rounds in 100\% Epon resin (each time for 8\,h or overnight) before embedding for curing at 60\,°C for 1–3\,d. All resin steps were carried out at 4\,°C. For each resin exchange, the tubes were taken out from the fridge 20–30\,min before for warming to RT.

\subsection*{Light microscopic imaging of the sample surface to assess resin gradient}
Resin-embedded samples (2-mm hemispheres) were trimmed to expose the center, and the surface was smoothed with a diamond knife. Light microscopic images of the sample surface were acquired with a slight tilt of the imaging plane, such that the otherwise black sample surface appeared silver colored.

\subsection*{X-ray microtomography volumetric imaging of samples without resin embedding}
To allow imaging by µCT without the need for resin embedding, 2-mm samples were embedded instead in 2\% agarose (Sigma-Aldrich) in 0.15\,M CaC or in water (depending on the last staining step that the sample was exposed to) in 2-ml Eppendorf tubes. They were imaged by μCT (Zeiss Xradia 520 Versa) using a voltage of 80\,kV at a voxel size of 3–6\,µm, using Zeiss Scout-and-Scan software (version 16.1.14271.44713). For µCT imaging of hemispheres, the samples were kept in the 50-ml glass tubes. Next, the glass tubes were put into a 140-ml syringe to be kept stable during µCT imaging (using a voxel size of 10–60\,µm). µCT datasets were visualized with Zeiss TXM3DViewer software (version 1.2.10).

\subsection*{Low-vacuum SEM imaging of incompletely stained samples}
To investigate samples at intermediate staining steps without complete staining (and therefore often reduced signal and conductivity), they were embedded in resin. Next, they were trimmed with a diamond head trimmer (Leica EM TRIM2) to expose the center to the block faces of samples; the block face was smoothed with a diamond knife ultramicrotome (Leica EM UC7) and then imaged in a scanning electron microscope with a field-emission cathode and low-vacuum mode (Quanta FEG 450, FEI). The chamber pressure was set to 30\,Pa. For the incident electron beam, a spot size of 3.5 at aperture 5 and an acceleration energy of 5\,keV (measured beam current, 67\,pA) were used for imaging at a pixel dwell time of 8–20\,μs and a pixel size of 11.24\,nm\textsuperscript{2} in plane (corresponding electron dose, 26–66\,e\textsuperscript{−}\,nm\textsuperscript{−2}) or 5.62\,nm\textsuperscript{2} in plane (corresponding electron dose, 105–264\,e\textsuperscript{−}\,nm\textsuperscript{−2}), at a working distance of about 5\,mm using the back scattered electron CBS detector.

\subsection*{High-vacuum SEM imaging of completely stained samples}
For those samples that were fully stained and embedded in resin (and therefore were expected to show sufficient conductivity), trimming and smoothing were similar to those of the previous section, but SEM imaging was performed in high-vacuum mode (5\,×\,10\textsuperscript{−4}\,Pa, Quanta FEG 420 and Quanta FEG 200). For the incident electron beam, a spot size of 3.2–3.5 and an acceleration energy of 2.8\,keV at aperture 4–5 (beam current, 72–160\,pA) were used. The pixel size was either 11.24\,nm (Quanta FEG 420; aperture 4; spot size, 3.5; pixel dwell time, 4–6\,µs; 31–47\,e\textsuperscript{−}\,nm\textsuperscript{−2}) (Quanta FEG 200; aperture 5; spot size, 3.2; dwell time, 8\,µs; 28\,e\textsuperscript{−}\,nm\textsuperscript{−2}) or 5.62\,nm (Quanta FEG 200; aperture 5; spot size, 3.2; dwell time, 8\,µs; 113\,e\textsuperscript{−}\,nm\textsuperscript{−2}).

\subsection*{Energy-dispersive X-ray spectroscopy (EDS) analysis}
Resin-embedded samples were trimmed to expose the center as a block face (Leica EM TRIM2), smoothed (Leica EM UC7) and coated with a 10-nm gold layer (Leica EM ACE600); afterward, they were imaged in a scanning electron microscope (Amray 1830) equipped with an Si(Li) EDS detector. An incident electron beam with an energy of 18\,keV was used at working distance of 15–20\,mm and a takeoff angle of 20.4°. The spectrum-collection time was 20–90\,s.

\subsection*{UV–visible spectrum acquisition of osmium(VI) solution}
We prepared 1\% potassium osmate(VI) in 0.15\,M sodium cacodylate buffer by adding 0.05\,g potassium osmate(VI) powder (Sigma-Aldrich) into 5\,ml cacodylate buffer. The solution was diluted 100 times and put into a glass cuvette to avoid spectrum signal clipping. The measurement was performed using the wavelength range 190–1,400\,nm on a UV–visible spectrometer (Jasco, V-670). Measurements on the same solution were made at time points 2\,min, 2\,h and 24\,h.

\subsection*{Raman spectrum measurement of osmium(VI) and FeCN solutions}
We measured the following chemicals without dilution in a custom-built Raman spectrometer (Chemistry Department of Goethe University Frankfurt) with a range of 0–4,400-cm\textsuperscript{−1} wavelength: (1) 0.15\,M \pdfpage{10}sodium cacodylate buffer; (2) 1\% potassium osmate(VI) in 0.15\,M cacodylate buffer; (3) 0.3\% potassium osmate(VI) and 1.3\% potassium ferrocyanide(II) (Sigma-Aldrich) in 0.15\,M cacodylate buffer; (4) 2.5\% potassium ferrocyanide(II) in 0.15\,M cacodylate buffer; (5) 1.9\% potassium ferricyanide(III) (Sigma-Aldrich) in 0.15\,M cacodylate buffer; (6) staining solution of 2-mm samples in 2-ml Eppendorf tubes, after 24\,h of 2\% OsO\textsubscript{4} (Serva) in 0.15\,M cacodylate buffer, 1\,h of 0.15\,M cacodylate buffer wash and 17\,h of 2.5\% potassium ferrocyanide(II) in 0.15\,M cacodylate buffer.

\subsection*{Membrane contrast quantification}
Membrane contrast was determined as the pixel-intensity differences between membrane and non-membrane regions in the EM images. For this, 500\,×\,500 pixel regions from the raw EM images (Fig. 1b and Extended Data Fig. 1m) were obtained and normalized for brightness (1st–99th percentiles, MATLAB 2018a, imadjust). Next, one randomly selected 500\,×\,500 subimage was used to train a random forest classifier (Fiji, Trainable Weka Segmentation\textsuperscript{26}) on ImageJ (version 1.53q) to automatically identify membrane versus non-membrane voxels. Next, this classifier was applied to all subimages, and membrane contrast was calculated as the difference between median membrane and non-membrane voxel intensities. To determine differences between the measured conditions, a two-way ANOVA was performed in GraphPad Prism 9, and Tukey’s multiple comparisons across staining protocols and sampling positions were computed. Next, to determine the significance of the observed differences, we performed one-sided \emph{t}-test for comparing the membrane contrast between the Hua protocol and our protocol (Os3 and Os24 pooled) at the core and intermediate positions, where improvements were expected.

\subsection*{Three-dimensional electron microscopy imaging and image alignment}
Subvolumes sized about (1.5\,mm)\textsuperscript{3} of the hemisphere samples H6 and H13 were cut out on the Leica EM TRIM2. SBEM datasets were acquired using a custom-built SBEM microtome\textsuperscript{27} mounted inside the chamber of a scanning electron microscope (FEI Quanta, Thermo Fisher Scientific). The image acquisition and SBEM microtome were controlled using custom-written software\textsuperscript{28}. An incident electron beam with an acceleration energy of 2.8\,keV at aperture 4 (spot size of 3.5; beam current, 160\,pA; dwell time, 2.8\,µs; electron dose, 22\,e\textsuperscript{−}\,nm\textsuperscript{−2}) or aperture 6 (spot size, 3.0; beam current, 47\,pA; dwell time, 12\,µs; electron dose, 28\,e\textsuperscript{−}\,nm\textsuperscript{−2}) was used to image at an in-plane pixel size of (11.24\,nm)\textsuperscript{2}.

The image alignment of all EM datasets and their segmentations were carried out with routines described in ref. 24. The tracing of axons was carried out with webKnossos\textsuperscript{29}. In brief, 20 axons were randomly seeded in a (3-µm)\textsuperscript{3} bounding box in the center of the H6 striatum dataset or from the control striatum dataset. Manual tracing was performed by three expert annotators. The occurrence of discontinuity (defined as not possible to trace non-ending nodes) was documented. The mean inter-error distance was calculated by dividing the total path length (mm) by the number of discontinuity nodes. Permutation random resampling was performed to calculate the 25th and 75th percentiles for each condition.

\subsection*{Statistics and reproducibility}
For the inter-error distance analysis (Fig. 2a), 20 randomly seeded axons from the dataset center were manually traced in the hemisphere and control striatum SBEM datasets, respectively. The mean inter-error distance was calculated as the total number of tracing errors divided by the total path length. Permutation (100 iterations) was performed to provide mean inter-error distance distributions for the hemisphere and control striatum datasets for the Kolmogorov–Smirnov test in MATLAB (kstest2).
\decompnote{As printed, this definition (errors divided by path length) is the inverse of the one in ``Three-dimensional electron microscopy imaging and image alignment'' (path length divided by the number of discontinuities) and of the plotted quantity in mm.}

For EDS analysis (Extended Data Figs. 1b,f and 4f), the elementary concentrations of Os, Fe and U (in wt \%) in the neuropil region from the center of a sample were measured. Data were presented as mean\,±\,s.d. as obtained from the EDS system. The staining experiments were replicated ≥3 times, while EDS was performed on one sample from each condition.

For the staining experiments on 2-mm and human samples, the number of replications was ≥3; for hemisphere and whole-brain staining experiments, the number of replications was ≥2 (Supplementary Table 3).

\emph{Note added in proof: A staining and embedding protocol for 2 mm-sized human cortical samples has recently been published (Karlupia et al., 2023,} https://doi.org/10.1016/j.biopsych.2023.01.025\emph{) with protocol steps partly overlapping with this protocols’ earlier version (Song et al., 2022, ref.}\,\textsuperscript{15}).\localcopy{karlupia2023}

\subsection*{Reporting summary}
Further information on research design is available in the Nature Portfolio Reporting Summary linked to this article.

\subsection*{Data availability}
All 3D EM datasets are publicly available for browsing at https://webknossos.org: H6 striatum (full-stained, full-volume embedded hemisphere, Extended Data Fig. 3), https://wklink.org/7781; H13 striatum (full-stained, subvolume embedded hemisphere, Extended Data Fig. 3), https://wklink.org/7718; H13 cortex (full-stained, subvolume embedded hemisphere, Extended Data Fig. 3), https://wklink.org/7723. The µCT datasets for full-embedded hemispheres and whole brains are also publicly available for browsing at https://webknossos.org: H6, https://wklink.org/7798; H4, https://wklink.org/7732 (note that this sample has one area of mechanical damage caused by sample handling in the peri-hippocampal region); W3, https://wklink.org/7742; W1, https://wklink.org/7789 (note that this sample was imaged before resin embedding, and it has one area of mechanical damage caused by sample handling in the frontal cortex). All source data have been deposited on the Edmond platform provided by the Max Planck Digital Library\textsuperscript{30}: https://doi.org/10.17617/3.RG58DU.

\subsection*{Code availability}
All analysis code (publicly available under the MIT license) has been deposited on the Edmond platform provided by the Max Planck Digital Library\textsuperscript{30}: https://doi.org/10.17617/3.RG58DU.

\subsection*{References}
\noindent 26.\ Arganda-Carreras, I. et al. Trainable Weka Segmentation: a machine learning tool for microscopy pixel classification. \emph{Bioinformatics} \textbf{33}, 2424–2426 (2017).\par
\noindent 27.\ Denk, W. \& Horstmann, H. Serial block-face scanning electron microscopy to reconstruct three-dimensional tissue nanostructure. \emph{PLoS Biol.} \textbf{2}, e329 (2004).\par
\noindent 28.\ Karimi, A. et al. Cell-type specific innervation of cortical pyramidal cells at their apical dendrites. \emph{eLife} \textbf{9}, e46876 (2020).\par
\noindent 29.\ Boergens, K. M. et al. webKnossos: efficient online 3D data annotation for connectomics. \emph{Nat. Methods} \textbf{14}, 691–694 (2017).\par
\noindent 30.\ Song, K. \emph{Source Data for ‘High-Contrast En-Bloc Staining of Mouse Whole-Brain Samples for EM-Based Connectomics’} https://doi.org/10.17617/3.RG58DU (2023).\par
\section*{Acknowledgements}
We thank Y. Hua for sharing very helpful insights into the staining protocols and for discussions and advice; E. Cocks for contributions to initial experiments and discussions; M. Sievers for initial contributions to resin protocols and comments on the manuscript; N. Ju for contributions to control experiments and human data; H.-S. Meyer and J. Gempt from the UKE \pdfpage{11}Department of Neurosurgery, Hamburg, Germany, for human sample access; A. Khalifa, Department of Connectomics, MPI for Brain Research, Frankfurt, for making the control striatum dataset available; and S. Loomba for support with SBEM data analysis. We are also grateful for brief discussions of chemical topics with D.M.E. van Niekerk, C. Sun, S. Chen, T.S. Hofer; to M. Koch for support with low-vacuum SEM; to M. Kind and T. Vavaleskou for support with EDS experiments; and to M. Braun, C. Boumrifak and Y. Aydogan for support with UV–visible spectrum and Raman spectrum measurements. We thank S. Soworka, S. Horn, L. Dadashev and I. Wolf for excellent technical support. All research was funded by the Max Planck Society.

\section*{Author contributions}
K.S. and M.H. conceived of the project. K.S., Z.F. and M.H. designed experiments. Z.F. carried out experiments with contributions from K.S.; K.S., M.H. and Z.F. analyzed data. K.S. and M.H. wrote the paper with contributions from Z.F.

\section*{Funding}
Open access funding provided by Max Planck Society.

\section*{Competing interests}
A patent application has been filed by the Max Planck Society with K.S., Z.F. and M.H. as inventors (application number EP22158620).

\section*{Additional information}
\noindent\textbf{Extended data} is available for this paper at https://doi.org/10.1038/s41592-023-01866-3.

\noindent \textbf{Supplementary information} The online version contains supplementary material available at https://doi.org/10.1038/s41592-023-01866-3.

\noindent \textbf{Correspondence and requests for materials }should be addressed to Kun Song or Moritz Helmstaedter.

\noindent \textbf{Peer review information} \emph{Nature Methods} thanks Martina Schifferer and the other, anonymous, reviewers for their contribution to the peer review of this work. Primary Handling Editor: Nina Vogt, in collaboration with the \emph{Nature Methods} team.
\decompnote{Left out here: the publisher's line ``Reprints and permissions information is available at www.nature.com/reprints.''}
 \pdfpage{12}\begin{center}
\includegraphics[width=0.95\linewidth,height=0.8\textheight,keepaspectratio]{figures/edfig1.png}
\end{center}

\textbf{Extended Data Fig. 1 | See next page for caption.}

\pdfpage{13}\textbf{Extended Data Fig. 1 | Development of 2\,mm protocol.} \textbf{(a)} Left: µCT image of staining gradient produced on 2\,mm samples with Hua protocol\textsuperscript{1} and confirmation with SEM (right). \textbf{(b)} EDS analysis of staining gradient from sample in (a). \textbf{(c)} µCT diagnosis of staining gradient after each consecutive step of the Hua staining protocol\textsuperscript{1}. \textbf{(d)} µCT images after extension of FeCN step from 3\,h to 12\,h. \textbf{(e)}. µCT images after extending TCH incubation. \textbf{(f)} Comparison of conductivity and membrane contrast of 1\,mm samples stained with TCH vs Pyrogallol. (left) samples stained by 1\,mm Hua protocol but with omission of TCH step; (middle) samples stained by 1\,mm protocol but replacing TCH by pyrogallol;\textsuperscript{2} (right) EDS analysis of the final osmium concentration in the corresponding samples. \textbf{(g)}. Contrast and gradient comparison of pyrogallol-related steps in water (top) vs. in CaC (bottom). Electron dose for high resolution EM images in f,g,i:18 e\textsuperscript{−}/nm\textsuperscript{2}\textbf{. (h)} Effect of duration of pyrogallol related steps. (left) influence of sampling location on the pyrogallol related gradient; (right) extension of pyrogallol incubation and 3\textsuperscript{rd} OsO\textsubscript{4} step removed gradient for all sampling locations. \textbf{(i)} Gradients related to UA and Ld incubation under different concentrations (UA only) and incubation time. \textbf{( j,k)} Resin infiltration speed in dependence of temperature and acetone concentration. Resin infiltration of larger samples requires separation of diffusion from polymerization to avoid softer core of sample (j).The resin fluidity time\textsuperscript{−}course curve (k) was measured for different conditions of temperature and acetone addition combination (see Methods). \textbf{(l)}. Quantification of membrane contrast (see Methods) in dependence of location within sample for Hua protocol (1b upper), our 2\,mm protocol-Os3 (Extended Data Fig. 1m) and Os24 (1b lower). \textbf{(m)} EM screening results for our 2 mm protocol-Os3.

\decompnote{As printed: the superscript citations ``Hua protocol\textsuperscript{1}'' and ``pyrogallol;\textsuperscript{2}'' in this caption do not match the article's reference list (ref.\ 1 is Karnovsky 1971, ref.\ 2 de Bruijn 1973; Hua et al.\ 2015 is ref.\ 6 and Mikula \& Denk 2015 ref.\ 13); they match the numbering of the Supplementary References (1 = Hua, 2 = Mikula). ``time\textsuperscript{−}course'' is printed with a superscript minus sign, as is ``whole\textsuperscript{−}brain'' in the caption of Extended Data Fig.\ 7.}
\alttext{Extended Data Fig.\ 1 fills p.\,12 (caption on p.\,13). Thirteen panels.
(a) µCT image (left, ``µCT'', scale bar) and SEM block-face image (right, ``SEM'', scale bar 500 µm) of a 2-mm sample: both show a bright rim and a darker, less stained core; in the SEM image a black arrow runs across the sample with dashed tick marks where the profile of (b) crosses the stain boundaries.
(b) Left: line plot, y axis ``Os intensity (a.u.)'' (−2 to 2, ticks every 1), x axis ``Location (µm)'' (0 to about 1,750, labels 0, 800, 1600), three curves ``µCT'' (red), ``SEM'' (blue), ``EDS'' (green) and four vertical dashed lines at about 340, 540, 1,110 and 1,410 µm. All three are high (about 1–2.5) in the outer 250 µm on both sides and low in the middle; the EDS curve drops to a flat floor at about −2.2 between about 560 and 1,100 µm, while µCT and SEM sit at about −1 to −1.5 there. Right: line plot with error bars, y axis ``WT.\%'' 0 to 30, x categories ``Periphery'', ``Middle'', ``Center''; Os (red) about 24, 26, 1.5; U (dark blue) about 22, 15.5, 0.5; Pb (green) about 8, 1, 0.3; Fe (grey) about 4.3, 1, 0.
(c) Row of five µCT images of 2-mm samples after the consecutive steps ``OsO\textsubscript{4} 1.5 h → FeCN1.5 h → OsO\textsubscript{4} 45 min → TCH 1 h → OsO\textsubscript{4} 1.5 h → UA'' (scale bar 1 mm): after the first OsO\textsubscript{4} step the sample is uniformly faint; after FeCN a faint rectangular ring appears; after the second OsO\textsubscript{4} a darker centre; after TCH and OsO\textsubscript{4} a bright rim with a dark core (two-layer gradient); after UA a bright sample with dark inner zones and cracks.
(d) ``OsO\textsubscript{4} 3 h-FeCN:'' µCT images after ``FeCN 3 h'', ``FeCN 7 h'', ``FeCN 12 h'' (scale bar 1 mm): a faint inner ring at 3 h that fades by 12 h.
(e) ``OsO\textsubscript{4} 3 h-FeCN 17-OsO\textsubscript{4} 3 h-TCH:'' µCT images ``TCH 3 h'' and ``TCH 5 h'' (scale bar 1 mm) with dark cracks and holes (breakages) that are more severe at 5 h, and a light-microscope photograph ``TCH 5 h'' (scale bar 2 mm) of a block face with many black holes.
(f) SEM images ``TCH replaced by H\textsubscript{2}O'' (inset overview with dark specks, scale bars 200 µm and 500 nm; the high-resolution image is a pale speckled field with barely visible membranes) and ``TCH replaced by Pg'' (scale bars 200 µm and 500 nm; clear dark membranes and mitochondria). Right: line plot ``wt.\%'' 0 to 60 for Os, Fe, U with error bars, ``TCH'' (blue) about 40, 1, 19; ``no TCH'' (orange) about 26, 1.5, 18; ``PG'' (red) about 37, 1.2, 16.
(g) SEM images ``H\textsubscript{2}O surround Pg'' and ``CaC surround Pg'', each with an overview inset (scale bar 500 µm) and a high-resolution image (scale bar 500 nm): with H\textsubscript{2}O the membranes are dark and sharp; with CaC the image is greyer with weaker membranes.
(h) Title ``OsO\textsubscript{4}-FeCN-Os-Pg-OsO\textsubscript{4}-UA-Ld''. Left: example µCT images ``no gradient'', ``weak gradient'', ``strong gradient'' (scale bar 1 mm; the strong-gradient sample has a dark centre) and a photograph of a resin disc with sampling positions 1, 2, 3 (scale bar 5 mm). Right: stacked bar chart ``ratio (\%)'' 0–100 for positions 1, 2, 3 at ``Pg: 6, 17, 17'' and ``OsO\textsubscript{4}: 3, 3, 6'' with colour key ``gradient: no'' (green), ``weak'' (yellow), ``strong'' (red). Pg 6/OsO\textsubscript{4} 3: position 1 about 62\% weak and 38\% no, position 2 about 66\% strong and 34\% weak, position 3 about 17\% strong, 17\% weak, 67\% no; Pg 17/OsO\textsubscript{4} 3: 100\% no, 66\% strong + 34\% no, 17\% weak + 83\% no; Pg 17/OsO\textsubscript{4} 6: 100\% no at all three positions.
(i) µCT images of four samples (scale bars 1 mm): ``2\% UA No Ld'', ``2\% UA Ld 50 °C, 4h'', ``4\% UA No Ld'' and ``4\% UA Ld 50 °C, 24h''; the samples with Ld show dark internal lens-shaped zones/cracks near the top.
(j) Sketch ``Resin infiltration vs. polymerization'' (an olive-green cylinder with wavy arrows entering from the sides) and light-microscope images of block faces ``100\%, RT'' (a darker, granular, softer-looking core) and ``95\%, 4°C'' (homogeneous; scale bar 500 µm).
(k) Line plot ``Resin fluidity (a.u.)'' 0 to 4 vs ``Time (days)'' 0 to 10: ``95\%, 4°C'' (red; four curves, three starting at about 3.3–3.5 between day 0 and day 0.7 and one starting at about 0.85 at day 4, all falling to about 0.05–0.8 by days 6–10), ``100\%, 4°C'' (light blue/violet; three curves starting at about 1.9 at day 0 and about 0.9 near day 1, below 0.3 by days 4–7) and ``100\%, RT'' (black lines with green dots; starting at about 1.6 and 1.2, reaching about 0 at days 3–4).
(l) Line plot with error bars, y axis ``Intensity difference membrane vs. background'' 0 to 100, x ``Periphery'', ``Intermediate'', ``Core'': ``This protocol'' blue (``Fig.1b bottom row'') about 91, 90, 93 and magenta (``m'') about 89, 88, 87; ``Hua (Fig.1b upper row)'' black about 85, 73, 64. Asterisks above Intermediate and Core.
(m) SEM ``Overview'' of a 2-mm sample (scale bar 500 µm, three white crosses) and high-resolution images ``Periphery'', ``Intermediate'', ``Core'' (scale bar 500 nm), all with similarly dark membranes.}
\figuredata{b (left)}{figures/edfig1bleft.csv}{figures/edfig1bleft.py}{source data (tier 1) for y: Edmond data/Extended Data Fig.1/Extended Data Fig. 1a-b/EDX/Suppl\_Fig1b.mat.mat; x in µm fitted to the raster figure (the source holds only sample indices); the curves labelled µCT and SEM in the figure are the .mat variables sem2n and uct2n, see CSV header}
\figuredata{b (right)}{figures/edfig1bright.csv}{figures/edfig1bright.py}{source data (tier 1): Edmond 'EDS raw data.doc' (EDAX reports Image1-1 to 1-3); printed points checked by digitizing the raster figure (they lie 0.2--0.5 wt.\% above the reports; printed bars are about half the 2-sigma errors)}
\figuredata{f (right)}{figures/edfig1fright.csv}{figures/edfig1fright.py}{source data (tier 1): Edmond 'Extended Data Fig. 1f/EDX/TCH.doc', 'no TCH.doc', 'PG.doc'; the plotted points are the second report of each file with ±2-sigma error, verified on the raster figure}
\figuredata{h (right)}{figures/edfig1h.csv}{figures/edfig1h.py}{digitized (tier 3) from the raster figure, ±1\%; no source data}
\figuredata{k}{figures/edfig1k.csv}{figures/edfig1k.py}{source data (tier 1): Edmond 'Extended Data Fig.1k/resin data.pzfx' (Prism XML); plotted columns identified on the raster figure}
\figuredata{l}{figures/edfig1l.csv}{figures/edfig1l.py}{source data (tier 1): Edmond 'Extended Data Fig. 1l/3. Calculated contrast values/contrast\_data.mat'; mean ± s.d. as in the authors' EDF\_Fig1l.m, checked on the raster figure (within 1 unit)} \pdfpage{14}\begin{center}
\includegraphics[width=0.95\linewidth,height=0.75\textheight,keepaspectratio]{figures/edfig2.png}
\end{center}

\textbf{Extended Data Fig. 2 | Effect of Os and FeCN incubation time and temperature } \textbf{on membrane contrast in 2\,mm samples.} \textbf{(a)} Membrane contrast comparison of samples stained with 3\,h of OsO\textsubscript{4} (left) or with 3\,h OsO\textsubscript{4} followed by 17\,h of FeCN. \textbf{(b)} Membrane contrast comparison of samples stained with 24\,h of OsO\textsubscript{4} (left), or with 24\,h of OsO\textsubscript{4} followed by 17\,h of FeCN. \textbf{(c,d)} Ultrastructural quality of samples stained with 3 (c) and 6 (d) days of OsO\textsubscript{4} at room temperature. \textbf{(e,f)} Ultrastructural quality of samples stained with 7d incubation of OsO\textsubscript{4} at 4\,°C (e) or with 6d incubation of OsO\textsubscript{4} at 4\,°C followed by 1 d incubation at RT (f). \textbf{(g-j)}. Membrane contrast and ultrastructural quality of samples stained with OsO\textsubscript{4} and FeCN under different incubation time and temperature: (g) 6d incubation of OsO\textsubscript{4} at 4\,°C followed by 1 d incubation at RT, plus 1d FeCN incubation at RT; (h) 6d incubation of OsO\textsubscript{4} at 4\,°C followed by 1d incubation at RT, plus 1d FeCN incubation at 4\,°C; (i) 7d with OsO\textsubscript{4} at 4\,°C followed by 1d of FeCN at RT; (j) 7d with OsO\textsubscript{4} at 4\,°C followed by 1d of FeCN at 4\,°C. All images were acquired in low vacuum SEM at 30\,Pa, with maximum electron dose 70 e\textsuperscript{−}/nm\textsuperscript{2}.

\alttext{Ten panels (a--j) of low-vacuum SEM images of 2-mm samples, each a pair: a lower-magnification image (scale bar 1 µm in c--j; the left images in a and b carry a scale bar without number) and a higher-magnification image (scale bar 500 nm). Each panel is titled with its staining condition, the varied part in red: (a) ``OsO\textsubscript{4} 3 h'' and ``OsO\textsubscript{4} 3 h --> FeCN 17 h'': both images grey and blurry, membranes and cytoplasm similarly stained, little membrane contrast. (b) ``OsO\textsubscript{4} 24 h'' and ``OsO\textsubscript{4} 24 h-->FeCN 17 h'': with FeCN the membranes are darker and sharper against a lighter cytoplasm. (c) ``OsO\textsubscript{4} 3 d RT'' and (d) ``OsO\textsubscript{4} 6 d RT'': high membrane contrast but large pale extracted areas around a cell body; in d a white box marks the zoomed region, which shows torn, extracted cytoplasm. (e) ``OsO\textsubscript{4} 7 d 4°C'' and (f) ``OsO\textsubscript{4} 6 d 4°C, 1 d RT'': well-preserved neuropil, moderately dark membranes, no extraction; f shows a few very dark myelin-like rings. (g) ``OsO\textsubscript{4} 6 d 4°C, 1d RT-->FeCN 1 d RT'': pale, fragmented, damaged tissue. (h) ``OsO\textsubscript{4} 6 d 4 °C, 1 d RT-->FeCN 1 d 4°C'': preserved tissue with dark membranes. (i) ``OsO\textsubscript{4} 7 d 4°C ->FeCN 1 d RT'' and (j) ``OsO\textsubscript{4} 7 d 4 °C ->FeCN 1 d 4°C'': both preserved with clear membranes and mitochondria.} \pdfpage{15}\begin{center}
\includegraphics[width=0.95\linewidth,height=0.75\textheight,keepaspectratio]{figures/edfig3.png}
\end{center}

\textbf{Extended Data Fig. 3 | Reproduction of mouse hemisphere stainings.} \textbf{(a-d)} Hemispheres H4-H6 (in total n\,=\,4 successful replicates) were stained as in Supplementary Fig. 7, Supplementary Table 3. \textbf{(e)} Whole brain W3 was stained as in Extended Data Fig. 7, Supplementary Table 3. Together with W1 (Fig. 1c), a total of n\,=\,2 whole brains were stained with parameters reported in Extended Data Fig. 7, Supplementary Table 3. All SEM images were taken at high vacuum with 5.62\,nm pixel size and a dose of 113 e\textsuperscript{−}/nm\textsuperscript{2}.

\decompnote{``Supplementary Fig. 7'' as printed; there is no Supplementary Fig.\ 7. The hemisphere protocol is shown in Extended Data Fig.\ 7.}
\alttext{A grid of five columns (a--e: hemispheres ``H4'', ``H5'', ``H6'', ``H13'' and whole brain ``W3'') and six rows. Row ``ov'': overview images of coronal block faces (scale bars 1 mm for H4--H13, 2 mm for W3) with small white crosses marking the sampling positions from the cortical surface to subcortical tissue. Rows ``L2/3'', ``L4'', ``L5'', ``WM'' and ``subcortical'': high-resolution SEM images (scale bar 500 nm in each, number printed in the bottom row), all with dark, continuous membranes, mitochondria and synaptic vesicles visible; the WM row shows dark myelin sheaths. Contrast is similar in all columns and rows.} \pdfpage{16}\begin{center}
\includegraphics[width=0.95\linewidth,height=0.8\textheight,keepaspectratio]{figures/edfig4.png}
\end{center}

\textbf{Extended Data Fig. 4 | See next page for caption.}

\pdfpage{17}\textbf{Extended Data Fig. 4 | Gradients and breakages in mouse hemispheres.} a Observation of OsO\textsubscript{4} diffusion into mouse hemisphere sample at successive time points at 4\,°C using µCT. \textbf{(b,c)} Effect of buffer rinsing between OsO\textsubscript{4} and FeCN incubation on staining gradient. \textbf{(d,e)} Effect of buffer rinsing between OsO\textsubscript{4} and FeCN incubation on FeCN diffusion speed and membrane contrast. \textbf{f} EDS measurement of Fe concentration in the center of samples from d,e. \textbf{g} gradient (top) after insufficient buffer washing step (1 or 2 days) at 4\,°C; gradient amplification and breakages occured in Pg (bottom). \textbf{h} No gradient after sufficient buffer washing (2 days RT or 4 days 4\,°C). \textbf{i} Effect of adding an extra OsO\textsubscript{4} incubation step after FeCN on the sample stability in H\textsubscript{2}O measured as sample expansion ratio. All SEM images were acquired in low vacuum SEM at 30\,Pa, with maximum electron dose 70 e\textsuperscript{−}/nm\textsuperscript{2}.

\alttext{Extended Data Fig.\ 4 fills p.\,16 (caption on p.\,17).
(a) ``OsO\textsubscript{4} 4°C'': sagittal µCT views of a hemisphere at ``17 h'', ``24 h'' and ``40 h'' (scale bar 2 mm), the bright stained rim growing inward with time; at right a line plot ``Depth (mm)'' 0 to 2.5 vs ``Incubation Time (h)'' 0 to 40: points at (0, 0), (17, 1.8), (24, 2.0), (40, 2.4) joined by lines.
(b) ``OsO\textsubscript{4}-FeCN'': coronal overview (scale bar 2 mm, three white crosses) and SEM images ``Periphery'', ``middle'', ``center'' (scale bar 500 nm); the periphery image is greyer, with less distinct membranes, than middle and center.
(c) ``OsO\textsubscript{4}-CaC-FeCN'': same layout (scale bar 2 mm, 500 nm); membranes similarly clear at all three positions.
(d) ``OsO\textsubscript{4}-FeCN1.5'': overview of a 2-mm sample (scale bar 500 µm) showing a bright outer ring and a darker centre, and SEM images ``Periphery'', ``middle'', ``center'' (scale bar 500 nm) with very dark membranes at the periphery and weaker ones at the center.
(e) ``OsO\textsubscript{4}-CaC-FeCN1.5'': homogeneous overview (scale bar 500 µm) and ``Periphery'' and ``center'' images (scale bar 500 nm) with similar contrast.
(f) Line plot ``Fe wt.\%'' 0 to 0.1 at ``Periphery'' and ``Center'': ``no CaC'' (dashed) about 0.085 ± 0.011 at the periphery falling to about 0.007 at the center; ``CaC wash'' (solid) about 0.047 at both positions, with error bars.
(g) Sagittal µCT views ``before Pg'' (top) and ``After Pg'' (bottom) of ``H1: cac 1 d 4 °C + pg-cac'' and ``H3: CaC 2 d 4 °C'' (scale bars 2 mm): before Pg a bright rim and darker interior; after Pg the H1 hemisphere is broken into pieces with dark interior regions and H3 shows a dark band and cracks.
(h) Sagittal µCT views ``After Pg'' and ``After 3rd Os'' of ``H4: CaC 2 d RT'' and ``H13: CaC 4 d 4 °C'' (scale bars 2 mm): intact and homogeneously bright.
(i) Photographs on mm grid paper (scale bars 1 mm) of two 2-mm sample pieces each: ``fresh'' (brown), and after ``OsO\textsubscript{4}-CaC-FeCN-CaC-(OsO\textsubscript{4} vs. no OsO\textsubscript{4})-CaC-H2O'' ``no OsO\textsubscript{4}'' (black, swollen, irregular) and ``OsO\textsubscript{4}'' (black, compact); right: line plot ``Expansion Ratio'' 1.0 to 2.0 vs ``Duration (h)'' 0 to 100: ``No OsO\textsubscript{4}'' (circles) about 1.33 at 5 h, 1.6 at 25 h, 1.79 at 97 h; ``With OsO\textsubscript{4}'' (squares) about 1.25–1.27 throughout.}
\figuredata{a (right)}{figures/edfig4a.csv}{figures/edfig4a.py}{digitized (tier 3) from the raster figure; no source data for this panel}
\figuredata{f}{figures/edfig4f.csv}{figures/edfig4f.py}{digitized (tier 3) from the raster figure; no source data for this panel}
\figuredata{i (right)}{figures/edfig4i.csv}{figures/edfig4i.py}{digitized (tier 3) from the raster figure; no source data for this panel} \pdfpage{18}\begin{center}
\includegraphics[width=0.95\linewidth,height=0.8\textheight,keepaspectratio]{figures/edfig5.png}
\end{center}

\textbf{Extended Data Fig. 5 | See next page for caption.}

\pdfpage{19}\textbf{Extended Data Fig. 5 | Experiments conducted for exploring possible chemical mechanisms underlying staining contrast generation.} \textbf{(a,b)} Membrane contrast after adding an intermediate CaC washing step between OsO\textsubscript{4} and FeCN. Imaged in SEM in high vacuum, electron dose 31 e\textsuperscript{−}/nm\textsuperscript{2}. \textbf{(c)} Membrane contrast after staining with potassium osmate (vi) (note yielded precipitation band). \textbf{(d)} UV-vis spectrum of potassium osmate (vi) in sodium cacodylate buffer over 24\,h. \textbf{(e,f)} 2\,mm samples stained with ‘reduced osmium’ protocol for 3\,h vs. 24\,h. \textbf{(g)} Membrane contrast after adding CaCl\textsubscript{2} into OsO\textsubscript{4} before FeCN incubation. \textbf{(h)} Membrane contrast after OsO\textsubscript{4} incubation at 4\,°C before FeCN incubation. SEM acquisition for panels c,e,f,g,h was performed in low vacuum at 30\,Pa, maximum electron dose of 70 e\textsuperscript{−}/nm\textsuperscript{2}. \textbf{(i-m)} Possible chemical logic underlying staining results \textbf{(i)} Comparison of bright-field color appearance of samples stained in OsO\textsubscript{4} for 2\,h, 24\,h, 24\,h (diluted to evaluate color), and stained in OsO\textsubscript{4}(24\,h)->CaC->FeCN (intermediate CaC wash to remove OsO\textsubscript{4}), top row. For comparison, pure staining solutions are shown below: Os(vi) and in vitro reaction of Os(vi) with FeCN (bottom). Note that sample solution after 24\,h OsO\textsubscript{4} incubation appears similarly colored as Os(vi) solution (pink); and sample solution after Os-Cac-FeCN staining appears similarly colored as Os(vi)+FeCN reaction in vitro (blue\textsuperscript{−}green). \textbf{( j)}. (left) Raman spectrum measurements of the product of Os(vi) + FeCN showing a shifted Os(vi) peak; (middle) in the wavelength range corresponding to FeCN(ii) and FeCN(iii), there was only signal consistent with FeCN(ii) peaks, but not FeCN(iii), suggesting the reaction between Os(vi) and FeCN was not a redox reaction. \textbf{(k)}. The same Os(vi) signal was observed in brain staining solution after FeCN staining and in-vitro Os(vi) + FeCN reaction. \textbf{(l)} Sketch summary of the possible chemical reactions (see in Supplementary Results). \textbf{(m)} Summary of Os(vi) coordination chemistry potentially relevant for staining mechanisms.

\alttext{Extended Data Fig.\ 5 fills p.\,18 (caption on p.\,19).
(a) ``OsO\textsubscript{4} 3 h-CaC 3 h-FeCN 17 h'' and (b) ``OsO\textsubscript{4} 24 h-CaC 3 h-FeCN 17 h'' (time labels in red): SEM images (scale bar 500 nm) with overview insets (scale bar 250 µm, a white cross at the imaging site); both show dark membranes and mitochondria.
(c) ``Os vi'': ``Periphery'' (inset overview, scale bar 250 µm; image scale bar 500 nm) with dark, granular, irregular staining, and ``Middle'' (scale bar 10 µm) with many black precipitate specks and pale cell bodies.
(d) Line plot ``Absorption (a.u.)'' 0 to 0.3 vs ``Wavelength (nm)'' 300 to 1000 (labels 400, 800): ``fresh'' (blue), ``2 h'' (red-orange), ``24 h'' (yellow) almost on top of each other: about 0.2 at 300 nm, a minimum of about 0.055–0.06 near 410 nm, a peak of about 0.095–0.105 near 510 nm and a slow decline to about 0.015–0.025 at 1000 nm. Inset photograph of two vials, one pink, one black.
(e) ``rO 3 h'' and (f) ``rO 24 h'' (red titles): SEM images (scale bar 500 nm) with round overview insets (scale bar 250 µm, white cross); at 3 h membranes are dark but the tissue looks extracted with pale patches; at 24 h membranes are darker against a lighter cytoplasm.
(g) ``OsO\textsubscript{4} 24 h, CaCl\textsubscript{2}, RT-->FeCN 17 h'' and (h) ``OsO\textsubscript{4} 24 h, 4°C-->FeCN 17 h'': lower-magnification images (scale bar 2 µm) with a white box and the zoomed region (scale bar 500 nm); membranes and cytoplasm are both grey, with low membrane contrast.
(i) Photographs of tubes: ``Samples stained'' with ``OsO\textsubscript{4}'' ``3 h'' (pale yellow solution), ``24 h'' (dark brown/black), ``24 h (dilute)'' (pink), and ``OsO\textsubscript{4} ->CaC ->FeCN'' (blue-green); ``Solution only'': ``Os (vi)'' (dark red/pink) and ``Os (vi) + FeCN'' (blue-green).
(j) Raman spectra, y axis ``Intensity (a.u.)'' 0 to 2 ×10\textsuperscript{4}, x axis ``Wavelength (cm\textsuperscript{-1})'' in two segments 750–1010 and 1900–2300: ``Os vi+FeCN'' (red) and ``Os vi'' (blue). Left: both curves share a broad band peaking at about 0.85 ×10\textsuperscript{4} near 830 cm\textsuperscript{-1}; the red curve has a sharp peak of about 1.6 ×10\textsuperscript{4} at about 880 cm\textsuperscript{-1} (dashed line), the blue one a peak of about 1.15 ×10\textsuperscript{4} at about 915 cm\textsuperscript{-1} (dashed line); a horizontal arrow labelled ``Os vi peak shift'' points from the blue to the red peak. Right: two red peaks of about 1.5 and 1.8 ×10\textsuperscript{4} near 2045 and 2085 cm\textsuperscript{-1} labelled ``FeCN peaks''; the blue curve stays near 0; an arrow ``No FeCNiii peak'' points at the flat region near 2140 cm\textsuperscript{-1}.
(k) Raman spectra 750–1010 cm\textsuperscript{-1}, y 0 to 2 (no multiplier printed): ``Osvi + FeCN'' (red) with a broad maximum of about 1.65 near 835 cm\textsuperscript{-1} and a peak of about 1.15 at about 880 cm\textsuperscript{-1}; ``Brain FeCN'' (green) with a plateau of about 1.2 between 800 and 840 cm\textsuperscript{-1} and a peak of about 1.1 at the same position, marked by a dashed line.
(l) Schematic in two rows of a cytoplasm/extracellular-matrix network (crossed lines) between two membranes (double lines). Key: ``Membrane'', ``Os(iv)'' (black dot), ``Os(vi)'' (magenta dot), ``Free Os(vi)'' (open magenta circle), ``Exposed Os(vi)'' (open magenta square), ``Cytoplasm / extracellular matrix'', ``FeCN'' (blue hexagon), ``Os(vi)-FeCN coordination'' (magenta circle or square with blue hexagons). Top row: ``OsO4 incubation (short)'' → black and magenta dots on the membranes and magenta dots plus free Os(vi) circles in the network → ``FeCN incubation'' (blue hexagons coordinate the free Os(vi)) → ``CaC washing'' → ``Later steps'': dots remain both on the membranes and in the network. Bottom row: ``OsO4 incubation (long)'' → exposed Os(vi) squares and free circles in the network → ``FeCN incubation'' (hexagons bind circles and squares) → ``CaC washing'' (the network is left without dots) → ``Later steps'': black dots on the membranes only.
(m) Diagram: ``Os vi'' connected by lines to four ligands, each with a structure and a note: ``hydroxide'' OH- (``weak coordination, easier dismutation''); ``CaC anion'' O=As(O\textsuperscript{-})(CH\textsubscript{3})\textsubscript{2} (``stronger coordination, hard to dismutate, slow diffusion''); ``formamide'' H–C(=O)–NH\textsubscript{2} (``even stronger coordination, harder to dismutate, faster diffusion''); ``cyanide'' CN\textsuperscript{-} (``strongest coordination, hardest to dismutate, fastest diffusion'').}
\figuredata{d}{figures/edfig5d.csv}{figures/edfig5d.py}{source data (tier 1): Edmond 'Extended Data Fig.5d/data.mat' (data4.os\_vi\_cac, \_2h, \_20h); agrees with the printed curves within 0.002 absorbance}
\figuredata{j}{figures/edfig5j.csv}{figures/edfig5j.py}{source data (tier 1): Edmond 'Extended Data Fig.6j,k/data\_batch2.mat', 'data\_batch3.mat', 'peak\_calibration\_allrounds.mat' (MATLAB tables decoded), processed as in the authors' Extended\_data\_Fig6jk.m; matches the printed 750--1010 cm\textsuperscript{-1} window; in 1900--2300 cm\textsuperscript{-1} the printed curves use a different baseline/scale (CSV header)}
\figuredata{k}{figures/edfig5k.csv}{figures/edfig5k.py}{source data (tier 1): Edmond 'Extended Data Fig.6j,k/data\_round4.mat'; the printed y scale (no multiplier, curves smoothed) was fitted to the raster figure, see CSV header} \pdfpage{20}\begin{center}
\includegraphics[width=0.95\linewidth,height=0.8\textheight,keepaspectratio]{figures/edfig6.png}
\end{center}

\textbf{Extended Data Fig. 6 | See next page for caption.}

\pdfpage{21}\textbf{Extended Data Fig. 6 | Summary of remaining artifacts in whole−brain staining protocol and alternative protocol approach for especially sensitive } \textbf{samples.} \textbf{(a)} μCT cross section of H5 showing locations with artifacts detailed in b-e. \textbf{(b)} Red arrow: ‘Vessel artifact’: disruption of blood vessel pericytes from surrounding neuropil. \textbf{(c)} Vessel artifact visible in μCT images of 2\,mm samples stained with OsO\textsubscript{4} RT 1 day or 4\,°C for 4 days. \textbf{(d,e)} Yellow arrows: Micro-breakages with width of less than 1 μm were occasionally observed in subcortical areas, in particular in highly myelinated regions. \textbf{(f)} During H\textsubscript{2}O incubation, the cerebellum was especially sensitive to macroscopic damage. \textbf{(g)} Changing the CaC rinsing step between 1\textsuperscript{st} OsO\textsubscript{4} and FeCN to 4\,°C improved the stability of the cerebellum for whole-brain staining. All SEM images were acquired at high vacuum, with maximum dose of 21 e\textsuperscript{−}/nm\textsuperscript{2}. \textbf{(h)}. Sketch of experiments performed to explore alternatives to incubation in plain H\textsubscript{2}O around Pg step, which may create high osmotic stress. \textbf{(i-l)} Influence of possible free Os (vi) in the sample on establishment of staining gradient during pyrogallol incubation (see Supplementary Results). \textbf{(m)} Samples from j further stained with Os in water resulted in acceptable membrane contrast. \textbf{(n)} Samples from j further stained with Os in CaC resulted in acceptable membrane contrast. \textbf{m,n} imaged under low-vacuum SEM (60 pa, beam current \textasciitilde{}55 pA), pixel size 11.24\,nm, electron dose \textasciitilde{}25 e\textsuperscript{−}/nm\textsuperscript{2}.

\alttext{Extended Data Fig.\ 6 fills p.\,20 (caption on p.\,21).
(a) µCT coronal section of hemisphere ``H5'' (scale bar 1 mm) with a red arrow at the cortex and yellow arrows at the subcortical white matter.
(b) Three SEM images (scale bars 10 µm, 10 µm, 2 µm) of blood vessels; red arrows point to vessel walls/pericytes detached from the neuropil, leaving pale gaps around the vessels.
(c) µCT images ``OsO\textsubscript{4} 24 h RT'' and ``OsO\textsubscript{4} 4 d 4ºC'' of 2-mm samples (scale bars 100 µm) with round overview insets (scale bars 1 mm); red arrows point to dark elongated vessel profiles.
(d) SEM images (scale bars 10 µm and 500 nm) with yellow arrows pointing to thin white micro-cracks through myelinated neuropil.
(e) SEM images (scale bars 50 µm and 5 µm) of fibre bundles with yellow arrows pointing to cracks.
(f) ``H5: after FeCN vs. in resin'': µCT/photograph pair of a cerebellum region (scale bar 5 mm) with white arrows (one labelled ``38.83'') at the damaged outer cerebellum.
(g) Images of ``W1 (CaC 4ºC)'' and ``W2 (CaC RT)'' brains in water (scale bars 5 mm); W2 is surrounded by many small dark fragments.
(h) Flow diagram ``Hemisphere'' → ``Staining Stage1'': ``1\textsuperscript{st}Os'' 72 h 4ºC → ``cac'' 96 h 4ºC → ``FeCN'' 72 h 4ºC → ``cac'' 48 h → black (stained) hemisphere → ``Staining Stage2'' with four alternatives: ``2\textsuperscript{nd} Os → CaC → Pg in h2o (i)''; ``CaC → CaC → Pg in h2o (j,m,n)'' (red); ``2\textsuperscript{nd} Os → KCl → Pg in h2o (k)''; ``2\textsuperscript{nd} Os → KCl → Pg in h2o (l)'' with ``brief pH adjust'' (red).
(i) ``CaC-Pg'', (j) ``CaC-Pg (no Os vi)'', (k) ``KCl-Pg'', (l) ``KCl-Pg (no Osvi)'': coronal µCT sections (scale bars 1 mm); i and k show a darker central region (gradient), j and l are homogeneous.
(m) ``j after Pg - CaC- Os h2o'' and (n) ``j after Pg - CaC-Os CaC'': µCT sections (scale bars 1 mm) and SEM images (scale bars 500 nm) with visible, moderately dark membranes and mitochondria.} \pdfpage{22}\begin{center}
\includegraphics[width=0.95\linewidth,height=0.75\textheight,keepaspectratio]{figures/edfig7.png}
\end{center}

\textbf{Extended Data Fig. 7 | Overview of 2\,mm, hemisphere and whole−brain protocol steps with incubation times and temperature steps visualized.} Illustration of the relevance of various staining and washing steps for the key challenges of whole\textsuperscript{−}brain staining: avoidance of gradients, peripheral tissue destruction, and macro- or microbreakages. Note for intermediate sized samples (larger than 2\,mm in any dimension), incubation times should be increased by roughly interpolating to the hemisphere protocol.

\alttext{A table-like chart. Columns: ``Incubation steps''; ``Incubation time (h)'' drawn as horizontal bars with numbers, key ``2 mm'' (grey bar), ``Hem.'' (black bar), ``W.B.'' (black dashed bar); a column coloured by ``Incubation temperature'' (colour bar 0°C blue to 60°C red: dark blue = 4°C, light blue = room temperature (RT), red = 50°C); a column of notes; and at the right sketches of brains (a dark core with pale rim, a grey hemisphere, a hemisphere broken into pieces) connected by arrows to the notes. Transcribed rows (step, solution: 2 mm / hemisphere / whole brain in h; temperature; note):
1a OsO\textsubscript{4}: – / 72 / 96 h; 2 mm n/a, Hem./W.B. 4°C; ``4°C critical to avoid OsO\textsubscript{4} over-exposure in periphery of large samples'' (arrow to a sketch with a dark centre).
1b OsO\textsubscript{4}: 24 / – / –; 2 mm RT, Hem./W.B. n/a.
2 Cac: 2 / 96 / 168 (the W.B. bar is drawn broken with ``...''); 4°C; ``4°C long enough critical, if RT: PSD lost, cerebellum damage; if too short, gradient in FeCN, breakages in Pg''.
3 FeCN: 17 / 72 / 72; 4°C; ``If step cancel: brekage in Pg. 4°C critical for large samples, if RT: brekages in FeCN''.
4 Cac: 2 / 48 / 96; RT; ``Critical for very large samples so free FeCN diffuse out. If cancel: gradient in 2nd Os \& breakages in Pg.''
5 OsO\textsubscript{4}: 3 / 48 / 48; RT; ``Critical additional OsO\textsubscript{4} step for the integrity in water steps of samples of all sizes.''
6 Cac: 1 / 24 / 72; RT; ``Important for free OsO\textsubscript{4} to diffuse out, to avoid OsO\textsubscript{4}-Pg interactions''.
7 H\textsubscript{2}O: 2 / 24 / 48; RT; ``Essential for enhanced contrast; Critical for removing residual Os vi to avoid Os vi-Pg interaction, otherwise gradients and breakages in Pg''.
8 Pyrogallol: 17 / 24 / 24; RT; ``Freshly prepared PG.''
9 H\textsubscript{2}O: 2 / 24 / 48; RT; ``Wash out remaining Pg to avoid Pg-OsO\textsubscript{4} interactions''.
10 OsO\textsubscript{4, H2O}: 6 / 48 / 96; RT; ``OsO\textsubscript{4} critical for amplify staining intensity \& sample conductivity''.
11 H\textsubscript{2}O: 1 / 24 / 48; RT; ``Wash out remaining OsO\textsubscript{4} to avoid OsO\textsubscript{4}-UA interactions''.
12a Uranyl Acetate: 17 / 48 / 48; 4°C; 12b Uranyl Acetate: 2 / 5 / 5; 50°C; note for 12a/12b: ``UA critical for sample conductivity. Note higher UA concentration (4\%) for sufficient diffufsion.''
13 H\textsubscript{2}O: 1 / 24 / 48; RT.
Lines from the notes of steps 2--7 lead to the sketches of a broken hemisphere and a hemisphere with a gradient. The same numbers are in tables/edfig7.csv.}
\decompnote{The quoted notes in the description keep the figure's spelling (``brekage'', ``diffufsion'').}


\end{document}
